% Cayley's Collected Papers, Volume 4, Pages 551-600 (Publication Pages 530-578)
% Paper [298]: Report on the Progress of the Solution of Certain Special Problems of Dynamics
% (Continued from previous pages)
\documentclass[11pt,a4paper]{article}
\usepackage[utf8]{inputenc}
\usepackage[T1]{fontenc}
\usepackage[english]{babel}
\usepackage{amsmath,amssymb,amsthm}
\usepackage{geometry}
\usepackage{hyperref}
\usepackage{fancyhdr}
\usepackage{titlesec}
\usepackage{enumitem}
\usepackage{array}

\geometry{margin=1in}

\pagestyle{fancy}
\fancyhf{}
\fancyhead[C]{\small Cayley, Collected Papers, Vol.\ IV, [298]}
\fancyfoot[C]{\thepage}
\renewcommand{\headrulewidth}{0.4pt}

\theoremstyle{definition}
\newtheorem*{remark}{Remark}

\title{\textbf{Paper [298]}\\[6pt]
\Large Report on the Progress of the Solution of\\
Certain Special Problems of Dynamics\\[12pt]
\large (Continuation: Publication Pages 530--578)}
\author{Arthur Cayley}
\date{}

\begin{document}
\maketitle
\thispagestyle{fancy}

\noindent\rule{\textwidth}{0.5pt}

\section*{Problem of Two Centres of Force (continued)}

\noindent\textsc{Article Nos.~52 to 64.}

\medskip

\textbf{52.} A case satisfying the required conditions is found to be
\[
\frac{x}{y} X \frac{y'}{y},
\]
or, what is the same thing,
\[
U = 2au + \frac{\beta}{u}, \qquad V = 2cv + \frac{\delta}{v},
\]
that is, besides the forces $\frac{\beta}{u^2}$, $\frac{\delta}{v^2}$, which vary as $(\text{dist.})^{-2}$, there are the forces $2au$, $2cv$, varying directly as the distance, and of the same amount at equal distances; or, what is the same thing, there is, besides the forces varying as $(\text{dist.})^{-2}$, a force varying directly as the distance, tending to a third centre midway between the other two, a case which is specially considered in the memoir; it is found that the functions in $r$, $\theta$ under the radicals (instead of rising only to the order 4) rise in this case to the order 6.

\textbf{53.} Among other cases are found the following, viz.:
\begin{enumerate}[label=\arabic*.]
    \item $U = \alpha u + \beta u^2 + \gamma u^4$,
    \item $U = \alpha u + \beta u^2$,
\end{enumerate}
\[
V = \delta v + \epsilon v^2 + \zeta v^4,
\]
where $\beta = \epsilon$, or else $\alpha\epsilon = \gamma\delta = 2\beta\epsilon$.

In regard to the subject of this second memoir of Lagrange, see post, Miscellaneous Problems, Liouville's Memoirs, Nos.~100 to 105.

\textbf{54.} In the \textit{M\'ecanique Analytique} (1st ed.~1788, and 2nd ed.~t.~II.~1813), Lagrange in effect reproduces his solution for the above-mentioned law of force (say $U = \frac{\alpha}{u} + 2\mu$, $V = \frac{\beta}{v} + 2\mu v$). There are even in the third edition a few trifling errors of work to be corrected. The remarks above referred to, as made by Lagrange in his first memoir, are also reproduced (see ante, Nos.~49 and 50).

\textbf{55.} Legendre, \textit{Exercices de Calcul Int\'egral}, t.~II.~(1817), and \textit{Th\'eorie des Fonctions Elliptiques}, t.~I.~(1825), uses $p$ and $q$ in the place of Euler's $p$, $q$; the forces are assumed to vary as $(\text{dist.})^{-2}$, and in consequence of the change Euler's cubic radicals are replaced by quartic radicals involving only even powers of $p$ and $q$ respectively; that is, the radicals are in a form adapted for the transformation to elliptic integrals; in certain cases, however, it becomes necessary to attribute to Legendre's variables $p$ and $q$ imaginary values. The various cases of the motion are elaborately discussed by means of the elliptic integrals; in particular Legendre notices certain cases in which the motion is oscillatory, and which, as he remarks, seem to furnish the first instance of the description by a free particle of only a finite portion of the curve which is analytically the orbit of the particle; there is, however, nothing surprising in this kind of motion, although its existence might easily not have been anticipated.

\textbf{56.} \S~26 of Jacobi's memoir ``Theoria Novi Multiplicatoris \&c.'' (1845) is entitled ``Motus puncti versus duo centra secundum legem Neutonianum attracti.'' The equations for the motion in space are by a general theorem given in the memoir ``De Motu puncti singularis'' (1842), reduced to the case of motion in a plane: viz. if $x$, $y$ are the coordinates, the centre point of the axis being the origin, and $y$ being at right angles to the axis, and if the distance of the centres is $2a$; then the only difference is that to the expression for $dt^2$ there is added a term $\frac{\alpha}{y^2}$, which arises from the rotation about the axis.

Two integrals are obtained, one the integral of Vis Viva, and the other of them an integral similar to one of those of Euler's or Lagrange's. And then $x'$, $y''$ being the differential coefficients of $x$, $y$ with regard to the time, the remaining equation may be taken to be $y'\,dx - x'\,dy = 0$, where $x'$, $y'$ are to be expressed as functions of $x$, $y$ by means of the two given integrals. This being so, the principle of the Ultimate Multiplier$^{(1)}$ furnishes a multiplier of this differential equation, and the integral is found to be
\[
\frac{y'\,dx - x'\,dy}{xy(x'^2 - y'^2) + a(a^2 - x^2 + y^2)xy} = d\psi,
\]
the quantity under the integral sign being a complete differential. To verify \textit{a posteriori} that this is so, Jacobi introduces the auxiliary quantities $\lambda'$, $\lambda''$ defined as the roots of the equation $\lambda^2 + \lambda(x^2 + y^2 - a^2) - a^2y^2 = 0$, which in fact, if as before $u$, $v$ are the distances from the centres, leads to
\[
u + v = 2\sqrt{a^2 + \lambda'}, \qquad u - v = 2\sqrt{a^2 + \lambda''},
\]
so that $\lambda'$, $\lambda''$ are functions of $u+v$, $u-v$ respectively; and the formul\ae, as ultimately expressed in terms of $\lambda'$, $\lambda''$, are substantially of the same form with those of Euler and Lagrange.

\footnotetext[1]{Explained in Jacobi's memoir ``Theoria Novi Multiplicatoris \&c.,'' Crelle, tt.~xxvii., xxviii., xxix.~1844--45.}

\textbf{57.} The investigations contained in Liouville's three memoirs ``Sur quelques cas particuliers \&c.'' (1846), find their chief application in the problem of two centres, and by leading in the most direct and natural manner to the general law of force for which the integration is possible, they not only give some important extension of the problem, but they in fact exhibit the problem itself and the preceding solutions of it in their true light. But as they do not relate to this problem exclusively, it will be convenient to consider them separately under the head Miscellaneous Problems.

\textbf{58.} In Serret's ``Th\`ese sur le Mouvement \&c.'' (1848), the problem is very elegantly worked out according to the principles of Liouville's memoirs as follows: viz. assuming that the expression of the distance between two consecutive positions of the body is
\[
ds^2 = \lambda\,(m\,d\mu^2 + n\,d\nu^2) + \lambda''\,d\psi^2,
\]
where $m$, $n$ are functions of $\mu$, $\nu$ respectively, and if the forces can be represented by means of a force-function $U$, then the motion can be determined, provided only $U$ be of the form
\[
U = \frac{\Phi\mu + \Psi\nu}{\lambda},
\]
where the functional symbols $\Phi$, $\Psi$, \&c.\ denote any arbitrary functions whatever.

\textbf{59.} It is then assumed that $\mu$, $\nu$ are the parameters of the confocal ellipses and hyperbolas situate in the moveable plane through the axis, viz. that we have
\[
\frac{x^2}{\mu^2} + \frac{y^2}{\mu^2 - b^2} = 1, \qquad \frac{x^2}{\nu^2} + \frac{y^2}{\nu^2 - b^2} = 1
\]
(the origin is midway between the two centres, $2b$ being their distance; $\mu$, $\nu$ are in fact equal to the sum and difference $u+v$, $u-v$ of the two centres respectively); and that the position of the moveable plane is determined by means of $\psi$, the inclination to a fixed plane through the axis, or say, as before, its azimuth. In fact, with these values of the coordinates, the expression of $ds^2$ is
\[
ds^2 = \frac{\mu^2 - \nu^2}{\mu^2 - b^2} \, d\mu^2 + \frac{\mu^2 - \nu^2}{b^2 - \nu^2} \, d\nu^2 + \mu^2\nu^2 \, d\psi^2,
\]
which is of the required form. And moreover if the forces to the two centres vary as $(\text{dist.})^{-2}$, and there is besides a force to the middle point varying as the distance, then
\[
U = \frac{\alpha}{\mu^2} + \frac{\beta}{\nu^2} + k(\mu^2 - \nu^2),
\]
whence (observing that $\lambda = \mu^2 - \nu^2$) $U$ is of the required form. The equations obtained by substituting for $U$ the above value give the ordinary solution of the problem.

\textbf{60.} Liouville's note to the last-mentioned memoir (1848) contains the demonstration of a theorem obtained by a different process in his second memoir, but which is in the present note, starting from Serret's formul\ae, demonstrated by the more simple method of the first memoir, viz., it is shown that the motion can be obtained if the two centres, instead of being fixed, revolve about the point midway between them in a circle in such manner that the diameter through the two centres always passes through the projection of the body on the plane of the circle. It will be observed that the circular motion of the two centres is neither a uniform nor a given motion, but that they are, as it were, carried along with the moving body.

\textbf{61.} In Desboves's memoir ``Sur le Mouvement d'un point mat\'eriel \&c.'' (1848), the author developes the solution of the foregoing problem of moving centres, chiefly by the aid of the method employed in Liouville's second memoir. And he shows also that the methods of Euler and Lagrange for the case of two fixed centres apply with modification to the more complicated problem of the moving centres.

\textbf{62.} The problem of two centres is considered in Bertrand's ``M\'emoire sur les \'Equations diff\'erentielles \&c.'' (1852), by means of Jacobi's form of the equations of motion, viz., the problem is reduced to a plane one by means of the addition of a force $\frac{\alpha}{y^2}$ (anti, No.~56).

\textbf{63.} Cayley's ``Note on Lagrange's Solution \&c.'' (1857) is merely a reproduction of the investigation in the \textit{M\'ecanique Analytique}; the object was partly to correct some slight errors of work, and partly to show what were the combinations of the differential equations, which give at once the integrals of the problem.

\textbf{64.} In \S~II. of Bertrand's ``M\'emoire sur quelques unes des formes \&c.'' (1857), the following question is considered, viz., assuming that the dynamical equations
\[
\frac{d^2x}{dt^2} = \frac{dU}{dx}, \qquad \frac{d^2y}{dt^2} = \frac{dU}{dy}
\]
have an integral of the form
\[
a = Px'^2 + Qx'y' + Ry'^2 + Sx' + Ty' + K
\]
(where $a$ is the arbitrary constant, and $P$, $Q$, \dots, $K$ are functions of $x$ and $y$), it is required to find the form of the force-function $U$. It is found that $U$ must satisfy a certain partial differential equation of the second order, the general solution of which is not known; but taking $U$ to be a function of the distance from any fixed point (or rather the sum of any number of such functions), it is shown that the only case in which the differential equations for the motion of a point attracted to a fixed centre of forces have an integral of the form in question is the above-mentioned one of two centres, each attracting according to the inverse square of the distance, and a third centre midway between them, attracting as the distance.


\section*{The Spherical Pendulum}
\noindent\textsc{Article Nos.~65 to 73.}

\medskip

\textbf{65.} The problem is obviously the same as that of a heavy particle on the surface of a sphere. I have not ascertained whether the problem was considered by Euler. Lagrange refers to a solution by Clairaut, Mem.\ de l'Acad., 1735.

The question was considered by Lagrange, \textit{Mec.\ Anal.}, 1st edit.\ p.~283. The angles which determine the position are $\psi$ the inclination of the string to the horizon, $\phi$ the inclination of the vertical plane through the string to a fixed vertical plane, or say the azimuth. And then forming the equations of motion, two integrals are at once obtained; these are the integrals of Vis Viva, and an integral of areas. And these give equations of the form $dt = \text{funct.}(\psi)\,d\psi$, $d\phi = \text{funct.}(\psi)\,d\psi$; so that $t$, $\phi$ are each of them given by a quadrature in terms of $\psi$, which is the point to which the solution is carried. It is noticed that $\psi$ may have a constant value, which is the case of the conical pendulum.

\textbf{66.} In the second edition, t.~XI.~p.~197 (1815), the solution is reproduced; only, what is obviously more convenient, the angles are taken to be $\psi$, the inclination to the vertical, $\phi$, the azimuth.

It is remarked that $\psi$ will always lie between a greatest value $\alpha$ and a least value $\beta$, and the integrals are transformed by introducing therein instead of $\psi$ the angle $\sigma$, which is such that
\[
\cos\psi = \cos\alpha \sin^2\sigma + \cos\beta \cos^2\sigma,
\]
by which substitution they assume a more elegant form, involving only the radical $\sqrt{1 + k^2(\cos\beta - \cos\alpha)\cos 2\sigma}$, where $k$ is a constant depending on $\cos\alpha$, $\cos\beta$; and the integration is effected approximately in the case where $\cos\beta - \cos\alpha$ is small.

Bravais has noticed, however, that by reason of some errors in the working out, Lagrange has arrived at an incorrect value for the angle $\Phi$, which is the apsidal angle, or difference of the azimuths for the inclinations $\alpha$ and $\beta$: see the 3rd edition (1855), Note VII., where M.~Bravais resumes the calculation, and he arrives at the value $\Phi = \frac{1}{2}\pi(1 + \frac{3}{8}\alpha\beta)$, $\alpha$ and $\beta$ being small.

Lagrange considers also the case where the motion takes place in a resisting medium, the resistance varying as velocity squared.

\textbf{67.} A similar solution to Lagrange's, not carried quite so far, is given in Poisson's \textit{M\'ecanique}, t.~I.~pp.~385 \textit{et seq.}~(2nd ed.~1833).

A short paper by Puiseux, ``Note sur le Mouvement d'un point mat\'eriel sur une sph\`ere'' (1842), shows merely that the angle $\Phi$ is $> \frac{1}{2}\pi$.

\textbf{68.} The ulterior development of the solution consists in the effectuation of the integrations by the elliptic and Jacobian functions. It is proper to remark that the dynamical problem the solution whereof by such functions was first fairly worked out is the more difficult one of the rotation of a solid body, as solved by Jacobi (1839), in completion of Rueb's solution (1834), post, Nos.~186 and 197.

\textbf{69.} In relation to the present problem we have Gudermann's memoir ``De pendulis sph\ae rici s \&c.'' (1849), who, however, does not arrive at the actual expressions of the coordinates in terms of the time; and the perusal of the memoir is rendered difficult by the author's peculiar notations for the elliptic functions$^{(1)}$.

\footnotetext[1]{The mere use of sn, cn, dn as an abbreviation of the somewhat cumbrous sinam, cosam, $\Delta$am of the \textit{Fundamenta Nova} is decidedly convenient.}

\textbf{70.} It would appear that a solution involving the Jacobian functions was obtained by Dur\`ege, in a memoir completed in 1849, but which is still unpublished; see \S~XX. of his \textit{Theorie der elliptischen Functionen} (1861), where the memoir is in part reproduced. It is referred to by Richelot in the Note presently mentioned.

\textbf{71.} We have next Tissot's \textit{Th\`ese de M\'ecanique}, 1852, where the expressions for the variables in terms of the time are completely obtained by means of the Jacobian functions H, $\Theta$, and which appears to be the earliest published one containing a complete solution and discussion of the problem.

\textbf{72.} Richelot, in the Note ``Bemerkungen zur Theorie des Raumpendels'' (1853), gives also, but without demonstration, the final expressions for the coordinates in terms of the time.

Donkin's memoir ``On a Class of Differential Equations \&c.'' (1855) contains (No.~59) an application to the case of the spherical pendulum.

\textbf{73.} The first part of the memoir by Dumas, ``Ueber die Bewegung des Raumpendels, \&c.'' (1855), comprises a very elegant solution of the problem of the spherical pendulum based upon Jacobi's theorem of the Principal Function (1837), and which is completely developed by the elliptic and Jacobian functions. The latter part of the memoir relates to the effect of the rotation of the Earth; and we thus arrive at the next division of the general subject.


\section*{Motion as affected by the Rotation of the Earth, and Relative Motion in general}
\noindent\textsc{Article Nos.~74 to 85.}

\medskip

\textbf{74.} Laplace (\textit{M\'ec.\ C\'eleste}, Book X.~c.~5) investigates the equations for the motion of a terrestrial body, taking account of the rotation of the Earth (and also of the resistance of the air), and he applies them to the determination of the deviations of falling bodies, \&c. He does not, however, apply them to the case of the pendulum.

\textbf{75.} We have also the memoir of Gauss, ``Fundamental-gleichungen \&c.'' (1804): the equations ultimately obtained are similar to those of Poisson. I have not had the opportunity of consulting this memoir.

\textbf{76.} Poisson, in the ``Memoire sur le mouvement des Projectiles \&c.'' (1838), also obtains the general equations of motion, viz. (omitting terms involving $n^2$), these may be taken to be
\begin{align*}
\frac{d^2x}{dt^2} &= 2n\sin\beta\,\frac{dy}{dt} + X, \\
\frac{d^2y}{dt^2} &= -2n\sin\beta\,\frac{dx}{dt} - 2n\cos\beta\,\frac{dz}{dt} + Y, \\
\frac{d^2z}{dt^2} &= 2n\cos\beta\,\frac{dy}{dt} + Z,
\end{align*}
(see p.~20), where the axes of $x$, $y$, $z$ are fixed on the Earth and moveable with it: viz., $z$ is in the direction of gravity; $x$, $y$ in the directions perpendicular to gravity, viz., $y$ in the plane of the meridian northwards, $x$ westwards; $g$ is the actual force of gravity as affected by the resolved part of the centrifugal force; $\beta$ is the latitude. There are some niceties of definition which are carefully given by Poisson, but which need not be noticed here.

\textbf{77.} Poisson applies his formul\ae\ incidentally to the motion of a pendulum, which he considers as vibrating in a plane; and after showing that the time of oscillation is not sensibly affected, he remarks that upon calculating the force perpendicular to the plane of oscillation, arising from the rotation of the Earth, it is found to be too small sensibly to displace the plane of oscillation or to have any appreciable influence on the motion---a conclusion which, as is well known, is erroneous. He considers also the motion of falling bodies, but the memoir relates principally to the theory of projectiles.

\textbf{78.} That the motion of the spherical pendulum is sensibly affected by the rotation of the Earth is the well-known discovery of Foucault; it appears by his paper, ``D\'emonstration Physique \&c.,'' \textit{Comptes Rendus}, t.~XXXII.~1851, that he was led to it by considering the case of a pendulum oscillating at the pole; the plane of oscillation, if actually fixed in space, will by the rotation of the Earth appear to rotate with the same velocity in the contrary direction; and he remarks that although the case of a different latitude is more complicated, yet the result of an apparent rotation of the plane of oscillation, diminishing to zero at the equator, may be obtained either from analytical or from mechanical and geometrical considerations. Some other Notes by Foucault on the subject are given, \textit{Comptes Rendus}, t.~XXXV.~(1853).

\textbf{79.} An analytical demonstration of the theorem was given by Binet, \textit{Comptes Rendus}, t.~XXXII.~(1851), and by Baehr (1853). Various short papers on the subject will be found in the \textit{Philosophical Magazine}, and elsewhere.

\textbf{80.} In regard to the above-mentioned problem of falling bodies, we have a Note by W.~S., \textit{Camb.\ and Dubl.\ Math.\ Journ.}\ t.~III.~(1848), containing some errors which are rectified in a subsequent paper, ``Remarks on the Deviation of Falling bodies, \&c.'' t.~IV.~(1849), by Dr Hart and Professor W.~Thomson.

\textbf{81.} The theory of relative motion is considered in a very general manner in M.~Quet's memoir, ``Des Mouvements relatifs en g\'en\'eral \&c.'' (1853). Suppose that $x$, $y$, $z$ are the coordinates of a particle in relation to a set of moveable axes; let $\xi$, $\eta$, $\zeta$ be the coordinates of the moveable origin in reference to a fixed set of axes, and treating the accelerations $\frac{d^2\xi}{dt^2}$, $\frac{d^2\eta}{dt^2}$, $\frac{d^2\zeta}{dt^2}$ as if they were coordinates, let these, when resolved along the moveable axes, give $u'$, $v'$, $w'$: suppose, moreover, that $p$, $q$, $r$ denote the angular velocities of the system of the moveable axes (or axes of $x$, $y$, $z$) round the axes of $x$, $y$, and $z$ respectively; $u'$, $v'$, $w'$, $p$, $q$, $r$ are considered as given functions of the time, and then, if
\begin{align*}
X &= \frac{d^2x}{dt^2} + 2\left(q\frac{dz}{dt} - r\frac{dy}{dt}\right) + x(q^2 + r^2) - y\left(pq - \frac{dr}{dt}\right) - z\left(pr + \frac{dq}{dt}\right) - u', \\
Y &= \frac{d^2y}{dt^2} + 2\left(r\frac{dx}{dt} - p\frac{dz}{dt}\right) + y(r^2 + p^2) - z\left(qr - \frac{dp}{dt}\right) - x\left(qp + \frac{dr}{dt}\right) - v', \\
Z &= \frac{d^2z}{dt^2} + 2\left(p\frac{dy}{dt} - q\frac{dx}{dt}\right) + z(p^2 + q^2) - x\left(rp - \frac{dq}{dt}\right) - y\left(rq + \frac{dp}{dt}\right) - w',
\end{align*}
it is shown that the equations of motion are to be obtained from the equation
\[
\Sigma m \left[(X - A')\,\delta x + (Y - B')\,\delta y + (Z - C')\,\delta z\right] = 0,
\]
where $\delta x$, $\delta y$, $\delta z$ are the virtual velocities of the particle $m$ in the directions of the moveable axes.

This equation is in fact obtained as a transformation of the equation
\[
\Sigma m \left[\left(\frac{d^2\xi}{dt^2} - A\right)\delta\xi + \left(\frac{d^2\eta}{dt^2} - B\right)\delta\eta + \left(\frac{d^2\zeta}{dt^2} - C\right)\delta\zeta\right] = 0,
\]
which belongs to a set of fixed axes of $\xi$, $\eta$, $\zeta$.

\textbf{82.} The equations for the motion of a free particle are of course $X = A'$, $Y = B'$, $Z = C'$.

In the case where the moveable axes are fixed on the Earth, and moveable with it (the diurnal motion being alone attended to), these lead to equations for the motion of a particle in reference to the Earth, similar to those obtained by Gauss and Poisson. The formul\ae\ are applied to the case of the spherical pendulum, which is developed with some care; and Foucault's theorem of the rotation of the plane of oscillation very readily presents itself. The general formulae are applied to the relative motion of a solid body, and in particular to the question of the gyroscope; the memoir contains other interesting results.

\textbf{83.} The principal memoirs on the motion of the spherical pendulum, as affected by the rotation of the Earth, are those of Hansen, ``Theorie der Pendelbewegung \&c.'' (1853), which contains an elaborate investigation of all the physical circumstances (resistance of the air, torsion of the string, \&c.) which can affect the actual motion, and the before-mentioned memoir by Dumas, ``Ueber der Bewegung des Raumpendels \&c.'' (1855).

The investigation is conducted by means of the variation of the constants; the integrals for the undisturbed problem were, as already noticed, obtained by means of Jacobi's Principal Function, that is, in a form which leads at once to the expressions for the variation of the constants; and the investigation appears to be carried out in a most elaborate and complete manner.

\textbf{84.} In concluding this part of the subject I refer to Mr Worm's work, \textit{The Rotation of the Earth} (1862), where the last-mentioned questions (falling bodies, the pendulum, and the gyroscope) are, in reference to the proofs they afford of the rotation of the Earth, considered as well in an experimental as in a mathematical point of view. The second part of the volume contains the theory (after Laplace and Gauss) of falling bodies, that of the pendulum (after Hansen), and that of the gyroscope (after Yvon Villarceau); and the whole appears to be a complete and satisfactory r\'esum\'e of the experimental and mathematical theories to which it relates.

\textbf{85.} We have also Cohen ``On the Differential Coefficients and Determinants of Lines \&c.'' (1862), where the equations for relative motion are obtained in a very elegant manner. The fundamental notion of the memoir may be considered to be the dealing directly with lines, velocities, \&c., which are variable in direction as well as in magnitude, instead of referring them, as in the ordinary analytical method, to axes fixed in space. The memoir is a highly interesting and valuable one, and the results are brought out with great facility; but I cannot but think that the great care required to apply the method correctly is an objection to it, if used otherwise than by way of interpretation of previously obtained results, and that the ordinary method is preferable.

I may remark that the theory of relative motion connects itself with the lunar and planetary theories as regards the reference of the plane of the orbit to the variable ecliptic, and as regards the variations of the position of the orbit; but this is a subject which I have abstained from entering upon.


\section*{Miscellaneous Problems}
\noindent\textsc{Articles Nos.~86 to 111 (several subheadings).}

\subsection*{Motion of a single particle}

\textbf{86.} Jacobi, in the memoir ``De Motu puncti singularis'' (1842), notices (\S~5) the case of a body acted on by a central force which is any homogeneous function of the degree $-2$ of the coordinates; or representing these by $r\cos\phi$, $r\sin\phi$, then the force is $= \frac{\Phi}{r^2}$, where $\Phi$ is any function of the angle $\phi$. The process is different from the ordinary one. In fact, after integrating by the process of a central force $= \frac{1}{r^2}$, he remarks that the method in fact applies to the more general law of force just mentioned.

\textbf{87.} Jacobi, in the memoir ``Theoria Novi Multiplicatoris \&c.'' (1845), considers (\S~25) the case of a body acted on by a central force $P$ a function of the distance, and besides by forces $X$ and $Y$, which are homogeneous functions of the degree $-3$ of the coordinates $(x, y)$; viz. the equations of motion are in this case
\[
\frac{d^2x}{dt^2} = \frac{Px}{r} + X, \qquad \frac{d^2y}{dt^2} = \frac{Py}{r} + Y,
\]
and there is an integral
\[
\tfrac{1}{2}(xy' - x'y)^2 - \int x^2(xY - yX)\,d\frac{y}{x} = \text{const.}
\]
(the function under the integral sign is obviously a function of the degree 0 in $(x, y)$, that is, it is a function of $\frac{y}{x}$).

If $X$, $Y$ are the derived functions of a force-function $U$ of the degree $-2$ in $(x, y)$, then there is, besides, the integral of Vis Viva, and thence a third integral is obtained by means of the theorem of the Ultimate Multiplier. It may be noticed that in the last-mentioned case the force-function is of the form $\frac{R}{r} + \frac{\Phi}{r^2}$, so that if we represent also the central force by means of a force-function $R$ ($= \text{function of } r$), then the entire force-function is $R + \frac{\Phi}{r^2}$. The case is a very interesting one; it includes that considered \S~IV.\ of Bertrand's ``Memoire sur les \'Equations diff\'erentielles de la M\'ecanique'' (1852), where the force-function is of the form $= \frac{\Phi}{r^2}$.

\subsection*{Motion of three mutually attracting bodies in a right line}

\textbf{88.} The problem is considered by Euler in the memoir ``De Motu rectilineo \&c.'' (1765), the forces being as the inverse square of the distance; and a solution is obtained for an interesting particular case. Let $A$, $B$, $C$ be the masses, and suppose that at the commencement of the motion the distances $CB$, $BA$ are in the ratio $\alpha : 1$, and that the velocities (assumed to be in the same sense) are proportional to the distances from a fixed point. Then, if $\alpha$ be the real root (there is only one) of the equation of the fifth order
\[
C(1 + 3\alpha + 3\alpha^2) = A\alpha^5(\alpha^2 + 3\alpha + 3) + B(\alpha + 1)^2(\alpha^2 - 1),
\]
the distances $CB$, $BA$ will always continue in the ratio $\alpha : 1$. It may be added that the distances $CB$, $BA$ each of them vary as $r^2 - a^2$, where $a$ is a constant, and $r$ is, according to the initial circumstances, a function of $t$ defined by one or the other of the two equations
\[
t = n\sqrt{r^2 - a^2} - na^2\log\frac{r + \sqrt{r^2 - a^2}}{a},
\]
or
\[
t = n\sqrt{a^2 - r^2} + na^2\sin^{-1}\frac{r}{a}.
\]

\textbf{89.} The bodies are considered as restricted to move in a given line; but it is clear that if the bodies, considered as free points in space, are initially in a line, and the initial velocities are also in this line, then the bodies will always continue in this line, which will be a fixed line in space. But if the distances and velocities are as above, except only that the velocities, instead of being along the line, are parallel to each other in any direction whatever, then the bodies will always continue in a line, which is in this case a moveable line in space (see post.\ No.~93).

\textbf{90.} Euler resumes the problem in the memoir of 1776 in the \textit{Nova Acta Petrop.} The distances $AB$, $BC$ being $p$ and $q$, then
\[
\frac{d^2p}{dt^2} = \frac{A + B}{p^2} + \frac{C}{q^2} - \frac{C}{(p+q)^2}, \qquad
\frac{d^2q}{dt^2} = \frac{A}{p^2} - \frac{A}{(p+q)^2} + \frac{B + C}{q^2},
\]
and in particular he considers the before-mentioned case of a solution of the form $p = \alpha q$; and also the particular problem where one of the masses vanishes, $C = 0$; in this case, introducing (instead of $p$, $q$), the new variables $u$, $s$, where $q = up$, $dq = s\,dp$ (a transformation suggested by the homogeneity of the equations), and making, moreover, the particular supposition that the integral of the first equation is $\left(\frac{dp}{dt}\right)^2 = \frac{2(A+B)}{p}$ (viz. making the constant of integration to vanish), he obtains between $s$ and $u$ the equation of the first order
\[
2(A + B)(s - u) = (A + B)s + A\frac{1}{(1+u)^2} - \frac{A}{u^2},
\]
which, however, he is not able to integrate.

\textbf{91.} Jacobi has given in the memoir ``Theoria Novi Multiplicatoris \&c.'' (1845) (\S~28, entitled ``De Problemate trium corporum in eadem recta, motorum. Substitutio Euleriana. Theoremata de viribus homogeneis'') a very symmetrical and elegant investigation of the same problem.

The centre of gravity being assumed to be at rest, the coordinates $x$, $x_1$, $x_2$ of the three bodies are in the first instance expressed as linear functions of the two variables $u$, $v$ (being, as Jacobi remarks, the transformation employed in his memoir ``Sur l'\'elimination des Noeuds'' (1843), post, No.~114); $\frac{d^2u}{dt^2}$ and $\frac{d^2v}{dt^2}$ come out respectively equal to homogeneous functions of the degree $-2$ of these variables $u$ and $v$, and the integral of Vis Viva exists.

The subsequent transformation consists in the introduction of the variables $r$, $\phi$, $s$, $\eta$, where $u = r\cos\phi$, $v = r\sin\phi$, $s = \frac{dr}{dt}$, $\eta = \sqrt{r^2}\frac{d\phi}{dt}$; this gives a system of equations independent of $r$; viz.,
\[
d\phi : ds : d\eta = \eta : \Phi s^2 + \eta^2 - \Psi : -\Phi s\eta + \Psi',
\]
where $\Phi$ is a given function of $\phi$, and $\Phi'$ is the derived function. If these equations were integrated, the equation of Vis Viva gives at once $r = \frac{1}{s}\sqrt{\Psi - \Phi(s^2 + \eta^2)}$; and finally the time $t$ would be given by a quadrature. The system of three equations has the multiplier $M = \frac{1}{s\eta - \Psi}$; hence if one integral were known the other would be at once furnished by the general theory. There is a simplification in the form of the solution if $h$ (the constant of Vis Viva) $= 0$. It is remarked that the method is equally applicable when the force varies as any power of the distance; and moreover that when the force varies as $(\text{dist.})^{-2}$, then the solution depends upon one quadrature only.

\textbf{92.} The concluding part of the section relates to the very general problem of a system of $n$ particles acted on by any forces homogeneous functions of the coordinates (this includes the case of $n$ particles mutually attracting each other according to a power of the distance), and this more general investigation illustrates the method employed in regard to the three bodies in a line. It may be remarked that in the general theorem for the $n$ particles ``sint vires \&c.,'' the constant of Vis Viva is supposed to vanish.

\subsection*{Particular cases of the motion of three bodies}

\textbf{93.} In the case of three bodies attracting each other according to the inverse square of the distance, the bodies may move in such manner as to be constantly in a line (a moveable line in space); this appears by the memoir, Euler, ``Consid\'erations g\'en\'erales \&c.'' (1764), in which memoir, however (which it will be observed precedes the memoir ``De Motu rectilineo \&c.'' (1765), referred to No.~88), Euler assumes that the mass of one of the bodies is so small as not to affect the relative motion of the other two. Calling the bodies the Sun, Earth, and Moon, and taking the masses to be $1$, $m$, $0$, then a result obtained is, that in order that the Moon may be perpetually in conjunction, its distance must be to that of the Sun as $a : 1$, where $m(1-a)^2 = 3a^5 - 3a^4 + a^3$, or $7 : 400$ nearly. It appears, however (anti, No.~88), that the foregoing restriction as to the masses is unnecessary, and, as will be mentioned, the problem has since been treated without such restriction. Euler investigates the motion in the case where the initial circumstances are nearly but not exactly as originally supposed; this assumes, however, that the motion is stable---i.e.\ that the bodies will continue to move nearly, but not exactly as originally supposed, which is at variance with the conclusions of Liouville's memoir, post, No.~95.

I have not examined the cause of this discrepancy.

\textbf{94.} In the \textit{M\'ecanique C\'eleste} (1799), Book X.~c.~6, Laplace considers two cases where the motion can be exactly determined.

$1^\circ$. Force varies as any function of the distance. It is shown that the motion may be such that the bodies form always an equilateral triangle of variable magnitude---the motion of each body about the centre of gravity being the same as if that point were a centre of force attracting the body according to a similar law.

$2^\circ$. Force $\propto (\text{dist.})^n$. The motion may be such that the three bodies are always in a right line (moveable in space), the relative distances being in fixed ratios to each other. In particular, if force $\propto (\text{dist.})^{-2}$, then $m$, $m'$, $m''$ being the masses, the quantity $z$ which determines the ratio of the distances $m''m'$, $mm$ is given by
\[
0 = mz^2[(1+z)^2 - 1] - m'(1+z)^2(1-z^2) - m''[(1+z)^2 - z^2] = 0,
\]
which is, in fact, the formula in Euler's memoir ``De Motu rectilineo \&c.''

\textbf{95.} Liouville's memoir ``Sur un cas particulier \&c.'' (1842) has for its object to show that if the initial circumstances are not precisely as supposed in the second of the two cases considered by Laplace, or, what is the same thing, in Euler's memoir ``Consid\'erations g\'en\'erales \&c.,'' then the motion is unstable; the instability manifests itself in the usual manner, viz. the expressions for the deviations from the normal positions are found to contain real exponentials which increase indefinitely with the time.

\textbf{96.} It may be proper to refer here to Jacobi's theorem, \textit{Comptes Rendus}, t.~III.~p.~61 (1836), quoted in the foot-note to No.~31 of my Report of 1857, [195], which relates to the motion of a point \textit{without mass} revolving round the Sun, and disturbed by a planet moving in a circular orbit, and properly belongs (as I have there remarked) to the problem of two centres, one of them moveable and the other revolving round it in a circle with uniform velocity. The theorem (given without demonstration by Jacobi) is proved by Liouville in his last-mentioned memoir, and he remarks that the theorem follows very simply as a corollary of the theorem by Coriolis, ``M\'emoire sur le principe des forces vives dans les mouvemens relatifs des Machines,'' \textit{Journ.\ de l'\'Ecole Polyt.}\ t.~XIII.~pp.~268---302 (1832).

There is, however, no difficulty in proving the theorem; another proof is given, Cayley, ``Note on a Theorem of Jacobi's \&c.'' (1862).

\subsection*{Motion in a resisting medium}

\textbf{97.} I do not consider the various integrable cases of the motion of a particle in a resisting medium, the resistance varying with the velocity according to some assumed law, the particle being either not acted on by any force, or acted upon by gravity only. Some interesting cases are considered in Jacobi's memoir ``De Motu puncti singularis'' (1842), \S~6 and 7 (see post, No.~108).

\textbf{98.} In the case of a central force varying as $(\text{dist.})^{-2}$, the effect of a resisting medium ($R \propto \text{vel}^2$) is considered in reference to the lunar theory, in the \textit{M\'ecanique C\'eleste}, Book VII.~c.~6. Formul\ae\ for the variations of the elliptic elements are given in the \textit{M\'ecanique Analytique}, t.~II.~(2nd edition). But the variations of the elliptic elements are fully worked out by means of elliptic and Jacobian functions in Sohncke's valuable memoir ``Motus Corporum \&c.'' (1833).

\textbf{99.} The effect of the resistance of the air on a pendulum has been elaborately considered by Poisson, Bessel, Stokes, and others; as the dimensions of the ball are attended to, the problem is in fact a hydrodynamical one. The effect on the spherical pendulum is considered in Hansen's memoir ``Theorie der Pendelbewegung \&c.'' (1853). The effect on the motion of a projectile is considered in Poisson's memoirs ``Sur le Mouvement des Projectiles \&c.'' (1838).

\subsection*{Liouville's memoirs ``Sur quelques Cas particuliers ou les \'equations du mouvement d'un point mat\'eriel peuvent s'int\'egrer'' (1846---49)}

\textbf{100.} In the first memoir (\S~1) the author considers a point moving in a plane or on a given surface, where the principle of Vis Viva holds good (or say where there is a force-function $V$). The coordinates of the point, and the function $U$, may be expressed in terms of two variables $\alpha$, $\beta$, and it is assumed that these are such that
\[
ds^2 = \lambda\,(d\alpha^2 + d\beta^2),
\]
where $\lambda$ is a function of $\alpha$ and $\beta$.

That is, we have $T = \frac{1}{2}\lambda(\alpha'^2 + \beta'^2)$; and the equations of motion are
\[
\frac{d}{dt}\lambda\alpha' = \frac{1}{2}\alpha'^2\frac{d\lambda}{d\alpha} + \frac{1}{2}\beta'^2\frac{d\lambda}{d\alpha} + \frac{dU}{d\alpha}, \qquad
\frac{d}{dt}\lambda\beta' = \frac{1}{2}\alpha'^2\frac{d\lambda}{d\beta} + \frac{1}{2}\beta'^2\frac{d\lambda}{d\beta} + \frac{dU}{d\beta}.
\]
One integral of these is $\lambda(\alpha'^2 + \beta'^2) = 2U + C$; and by means of it the equations take the form
\[
\frac{d}{dt}\lambda\alpha' = \beta'^2\frac{d\lambda}{d\alpha} + \frac{dU}{d\alpha}, \qquad
\frac{d}{dt}\lambda\beta' = \alpha'^2\frac{d\lambda}{d\beta} + \frac{dU}{d\beta}.
\]
These equations, it is easy to show, may be integrated if
\[
(2U + C)\lambda = f\alpha - F\beta,
\]
and they then in fact give
\[
\lambda\alpha' = \sqrt{2f\alpha + C\phi\alpha - A}, \qquad
\lambda\beta' = \sqrt{A - 2F\beta - C\phi\beta},
\]
where $A$ is an arbitrary constant.

And we then have
\[
\frac{d\alpha}{\sqrt{2f\alpha + C\phi\alpha - A}} = \frac{d\beta}{\sqrt{A - 2F\beta - C\phi\beta}},
\]
which gives the path, and
\[
\frac{ds}{\sqrt{2U + C}} = \frac{\sqrt{\lambda}\,d\alpha}{\sqrt{2f\alpha + C\phi\alpha - A}} = \frac{\sqrt{\lambda}\,d\beta}{\sqrt{A - 2F\beta - C\phi\beta}},
\]
the expression for the time is easily obtained by means of a quadrature.

It is not more general, but it is frequently convenient to employ instead of $\alpha$, $\beta$, two variables $\mu$ and $\nu$, such that
\[
ds^2 = \lambda\,(m\,d\mu^2 + n\,d\nu^2),
\]
where $m$ is a function of $\mu$ only and $n$ of $\nu$ only, while $\lambda$ contains $\mu$ and $\nu$. The geometrical signification of the equation $ds^2 = \lambda\,(d\alpha^2 + d\beta^2)$, or of the last-mentioned equivalent form, is that the curves $\alpha$ or $\mu = \text{const.}$, $\beta$ or $\nu = \text{const.}$ intersect at right angles.

The foregoing differential equation of the path, writing $f\mu$, $F\nu$ in the place of $f\alpha$, $F\beta$ respectively, may be expressed in the form
\[
f\mu \cos^2 i + F\nu \sin^2 i = A,
\]
where $i$, $90^\circ - i$ are the inclinations of the path at the point $(\lambda, \mu)$ to the two orthotomic curves through this point.

\textbf{101.} The before-mentioned equation $(2U + C)\lambda = f\alpha - F\beta$ may be satisfied independently of $C$, or else only for a particular value of $C$. In the former case the law of force is much more restricted, but on the other hand there is no restriction as regards the initial circumstances of the motion; it is the more important one, and is alone attended to in the sequel of the memoir. In the case in question (changing the functional symbols) we must have
\[
\lambda = \phi\alpha - \psi\beta, \qquad U = f\alpha - F\beta;
\]
so that the functions denoted above by $f\alpha$, $F\beta$ are $2f\alpha + C\phi\alpha$, $2F\beta + C\psi\beta$; the equation of the trajectory is
\[
\frac{d\alpha}{\sqrt{2f\alpha + C\phi\alpha - A}} = \frac{d\beta}{\sqrt{A - 2F\beta - C\psi\beta}},
\]
and for the time the formula is
\[
t = \int \frac{\phi\alpha\,d\alpha}{\sqrt{2f\alpha + C\phi\alpha - A}} + \int \frac{\psi\beta\,d\beta}{\sqrt{A - 2F\beta - C\psi\beta}}.
\]
It is noticed also that taking $B$, $\epsilon$ to denote two new arbitrary constants, and writing
\[
\int \frac{d\alpha}{\sqrt{2f\alpha + C\phi\alpha - A}} + \int \frac{d\beta}{\sqrt{A - 2F\beta - C\psi\beta}} = B,
\]
the equation of the trajectory and the expression for the time assume the forms
\[
t = \epsilon + \frac{dB}{dC}, \qquad 0 = \frac{dB}{dA},
\]
as is known \textit{a priori} by a theorem of Jacobi's.

If the forces vanish, the path is a geodesic line; and denoting by $a$ the ratio of the constants $A$, $C$, we have
\[
\frac{d\alpha}{\sqrt{\phi\alpha - a}} = \frac{d\beta}{\sqrt{a - \psi\beta}},
\]
and moreover
\[
ds = d\alpha\,\sqrt{\phi\alpha - a} + d\beta\,\sqrt{a - \psi\beta},
\]
which are geometrical properties relating to the geodesic line.

\textbf{102.} Passing to the applications: in the first place, if $\alpha$, $\beta$ are rectangular coordinates of a point in piano, then writing instead of them $x$, $y$, we have $ds^2 = dx^2 + dy^2$, which is of the required form; but the result obtained is the self-evident one, that the equations may be integrated by quadratures when $U$ is of the form $\text{funct.}\,x - \text{funct.}\,y$. But taking instead the elliptic coordinates $\mu$, $\nu$, of a point in piano---viz., as employed by the author, these are the semiaxes of the confocal ellipse and hyperbola represented by the equations
\[
\frac{x^2}{\mu^2} + \frac{y^2}{\mu^2 - b^2} = 1, \qquad \frac{x^2}{\nu^2} + \frac{y^2}{\nu^2 - b^2} = 1
\]
---very interesting results are obtained.

The equations give
\[
\delta x^2 = \mu^2\nu^2, \qquad \delta y^2 = (\mu^2 - b^2)(b^2 - \nu^2),
\]
and thence
\[
ds^2 = \frac{\mu^2 - \nu^2}{\mu^2 - b^2}\,d\mu^2 + \frac{\mu^2 - \nu^2}{b^2 - \nu^2}\,d\nu^2,
\]
which is of the proper form, and the corresponding expression of $\lambda U$ is
\[
\lambda U = f\mu - F\nu,
\]
so that the force-function having this value ($f\mu$, $F\nu$ being arbitrary functions of $\mu^2 - \nu^2$ and $\nu^2 - \mu^2$ respectively), the equations of motion may be integrated by quadratures.

\textbf{103.} In particular, if
\[
f\mu = g\mu + g'\mu^2 + k(\mu^4 - b^2\mu^2), \qquad F\nu = g\nu - g'\nu^2 - k(\nu^4 - b^2\nu^2),
\]
then
\[
U = \frac{g}{\mu + \nu} + \frac{g'}{\mu - \nu} + k(\mu^2 + \nu^2 - b^2).
\]
But $\mu + \nu$, $\mu - \nu$ are the distances of the point from the two foci, and $\mu^2 + \nu^2 - b^2$ ($= x^2 + y^2$) is the square of the distance from the centre, so that the expression for $U$ is
\[
U = \frac{g}{r} + \frac{g'}{r'} + kR^2;
\]
and the case is that of forces to the foci varying inversely as the squares of the distances, and a force to the centre varying directly as the distance---the case considered by Lagrange in the problem of two centres. But this is merely one particular case of those given by the general formula.

The cases $g = 0$, $g' = 0$, $k = 0$ (no forces), and $g = 0$, $g' = 0$ (a force to the centre) lead to some interesting results; it is noticed also that the expression for the force-function may be written
\[
U = \text{funct.}(r + r') \frac{1}{r + r'} + \text{funct.}(r - r') \frac{1}{r - r'},
\]
and that it may be ascertained (without transforming to elliptic coordinates) whether a given value of the force-function is of the form considered in the theory.

In \S~3 the author considers the expression $dx^2 + dy^2 = \lambda\,(d\alpha^2 + d\beta^2)$, $\lambda$ being in the first instance any function whatever of $\alpha$ and $\beta$; and he shows that the expressions of $x$, $y$ are given by the equation
\[
x + y\sqrt{-1} = \psi(\alpha + \beta\sqrt{-1}),
\]
$\psi$ being any real function. If, however, it is besides assumed that $\lambda$ is of the required form $= f\alpha - F\beta$, then he shows that the system of elliptic coordinates is the only one for which the conditions are satisfied. \S\S~4, 5, 6 and 7 relate to the motion of a point on a sphere, an ellipsoid, a surface of revolution, and the skew helicoid respectively; and the concluding \S~8 contains only a brief reference to the author's second memoir.

\textbf{104.} Liouville's second and third memoirs may be more briefly noticed. In the second memoir the author starts from Jacobi's theorem of the $V$ function, viz., assuming that there is a force-function $U$ independent of the time, then in order to integrate the equations of motion
\[
\frac{d^2x}{dt^2} = \frac{dU}{dx}, \qquad \frac{d^2y}{dt^2} = \frac{dU}{dy}, \qquad \frac{d^2z}{dt^2} = \frac{dU}{dz},
\]
it is required to find a function $\Theta$ of $x$, $y$, $z$ containing three arbitrary constants $A$, $B$, $C$ (distinct from the constant attached to $\Theta$ by mere addition) satisfying the differential equation
\[
\left(\frac{d\Theta}{dx}\right)^2 + \left(\frac{d\Theta}{dy}\right)^2 + \left(\frac{d\Theta}{dz}\right)^2 = 2U + k,
\]
for then the required integrals of the equations of motion are
\[
\frac{d\Theta}{dA} = A', \qquad \frac{d\Theta}{dB} = B', \qquad \frac{d\Theta}{dC} = C', \qquad t = \frac{d\Theta}{dk} + t',
\]
$A'$, $B'$, $C'$ being new arbitrary constants.

Liouville introduces in place of $x$, $y$, $z$, the elliptic coordinates $\rho$, $\mu$, $\nu$, which are such that
\[
\frac{x^2}{a^2 + \rho} + \frac{y^2}{b^2 + \rho} + \frac{z^2}{c^2 + \rho} = 1, \quad
\frac{x^2}{a^2 + \mu} + \frac{y^2}{b^2 + \mu} + \frac{z^2}{c^2 + \mu} = 1, \quad
\frac{x^2}{a^2 + \nu} + \frac{y^2}{b^2 + \nu} + \frac{z^2}{c^2 + \nu} = 1,
\]
or, what is the same thing,
\[
x^2 = \frac{(a^2 + \rho)(a^2 + \mu)(a^2 + \nu)}{(b^2 - a^2)(c^2 - a^2)}, \quad
y^2 = \frac{(b^2 + \rho)(b^2 + \mu)(b^2 + \nu)}{(c^2 - b^2)(a^2 - b^2)}, \quad
z^2 = \frac{(c^2 + \rho)(c^2 + \mu)(c^2 + \nu)}{(a^2 - c^2)(b^2 - c^2)},
\]
and he then finds that the resulting partial differential equation in $\rho$, $\mu$, $\nu$ may be integrated provided that $U$ is of the form
\[
U = \frac{f\rho}{(\rho - \mu)(\rho - \nu)} + \frac{F\mu}{(\mu - \nu)(\mu - \rho)} + \frac{\omega\nu}{(\nu - \rho)(\nu - \mu)},
\]
$f$, $F$, $\omega$ being any functional symbols whatever: viz., the expression for $\Theta$ is
\[
\Theta = \int\!\frac{d\rho}{\sqrt{2f\rho + A + B\rho^2 + 2C\rho^3}}
+ \int\!\frac{d\mu}{\sqrt{2F\mu + A + B\mu^2 + 2C\mu^3}}
+ \int\!\frac{d\nu}{\sqrt{2\omega\nu + A + B\nu^2 + 2C\nu^3}}.
\]
In the case where $U = 0$ we have a particle not acted on by any forces, and the path is of course a straight line. The peculiar form in which these equations are obtained leads to very interesting results in regard to the theory of Abelian integrals, and to that of the geodesic lines of an ellipsoid.

The formul\ae\ require to be modified in certain cases, such as $c = b$ or $b = 0$. The case $b = 0$ leads to the theory developed in the first memoir in relation to the problem of two centres. The case is indicated where $b = 0$, $c = 0$, the ratio $b : c$ remaining finite. The case is briefly considered of a particle moving on a given surface.

\textbf{105.} The third memoir purports to relate to a system of particles, but the formulae are exhibited under a purely analytical point of view; so much so, that the coordinates of the points (3 for each point) are considered as forming a single system of variables $x_1$, $x_2$, \dots, $x_n$. The partial differential equation is
\[
\left(\frac{d\Theta}{dx_1}\right)^2 + \left(\frac{d\Theta}{dx_2}\right)^2 + \cdots + \left(\frac{d\Theta}{dx_n}\right)^2 = 2U + k,
\]
which is transformed by introducing therein the new variables $\mu_1$, $\mu_2$, \dots, $\mu_n$ analogous to the elliptic coordinates of the second memoir. The memoir really belongs rather to the theory of the Abelian integrals (in regard to which it appears to be a very valuable one) than to dynamics.


\subsection*{Memoirs by Jacobi, Bertrand, and Donkin, relating to various Special Problems}

\textbf{106.} I have inserted this heading for the sake of showing at a single view what are the special problems incidentally considered in the undermentioned memoirs which are referred to in several places in the present Report.

\textbf{107.} Jacobi, ``De Motu puncti singularis'' (1842).---I call to mind that the memoir chiefly depends on the theorem of the Ultimate Multiplier (the theory in its generality being developed in the later memoir ``Theoria Novi Multiplicatoris \&c.,'' 1844---45). \S~4 is entitled ``The motion of a point on the surface of revolution,'' which, the principle of the conservation of areas holding good, is reduced to the problem of the motion on the meridian curve, and thus depends upon quadratures only. \S~5 is entitled ``On the motion of a point about a fixed centre attracted according to a certain law more general than the Newtonian one'' (anti, No.~86). \S~6, ``On the motion of a point on a given curve and in a resisting medium'' (resistance $= a + bv^2$, or $= a + bv^3$); and \S~7, ``On the Ballistic Curve,'' viz., the forces are gravity and a resistance $= a + bv^2$.

\textbf{108.} In Jacobi's memoir ``Theoria Novi Multiplicatoris \&c.'' (1845), \S~25 is entitled ``On the motion of a point attracted towards a fixed centre'' (see ante, No~87); \S~26, ``On the motion of a point attracted towards two fixed centres according to the Newtonian law'' (ante, No.~56); \S~27, ``On the rotation of a solid body about a fixed point'' (post.\ No.~193); \S~28, ``On the problem of three bodies moving in a right line; the Eulerian substitution; theorems on homogeneous forces'' (ante.\ No.~91); and \S~29, ``The principle of the ultimate multiplier applied to a free system of material points moving in a resisting medium; on the motion of a comet in a resisting medium about the sun.''

\textbf{109.} And in Jacobi's memoir ``Nova Methodus \&c.'' (1802), besides \S~64 and \S~65, which are applications of the method to general dynamical theorems, we have \S~66, containing a simultaneous solution of the problem of the motion of a point attracted to a fixed centre and of that of the rotation of a solid body (post.\ No.~206) and \S~67, relating to the motion of a point attracted to a fixed centre according to the Newtonian law.

\textbf{110.} Bertrand's ``M\'emoire sur les int\'egrales diff\'erentielles de la M\'ecanique'' (1852).---\S~III.\ relates to the motion of a point attracted to a fixed centre by a force varying as a function of the distance; \S~IV.\ to the case where the forces arise from a force-function $U = \frac{\Phi(\frac{y}{x})}{r^2}$ [or, what is the same thing, $= \frac{\Phi}{r^2}$] (anth.\ No.~87); \S~V.\ to the problem of two centres (anth.\ No.~62), and \S~VI.\ to the problem of three bodies (post.\ No.~117).

\textbf{111.} Donkin's memoir ``On a Class of Differential Equations \&c.'' (1855). Part I.\ Nos.~27 to 30 relate to the problem of central forces (in space), No.~31 to the rotation of a solid body, and \S~III.\ to the same subject, viz.\ Nos.~40 and 41 to the general case, Nos.~42 to 44 to the particular case $A = B$; and Nos.~45 to 48 to the reduction thereto of the general case by treating the forces which arise from the inequality of $A$ and $B$ as disturbing forces. Part II.\ Nos.~59 and 60 relate to the spherical pendulum; Nos.~72 and 73 to ``Transformation from fixed to moving axes of coordinates,'' say to Relative Motion; and Nos.~84 to 96 to the problem of three bodies (post.\ No.~120).


\section*{The Problem of Three Bodies}
\noindent\textsc{Article Nos.~112 to 123.}

\medskip

\textbf{112.} A system of differential equations, such as
\[
\frac{dx_1}{X_1} = \frac{dx_2}{X_2} = \cdots = \frac{dx_{n+1}}{X_{n+1}}
\]
($n$ equations between $n+1$ variables), may be termed a system of the $n$th order, or more simply a system of $n$ equations. Let $(u_1, u_2, \dots, u_{n+1})$ be any functions of the original variables $(x_1, x_2, \dots, x_{n+1})$, the system may be transformed into the similar system
\[
\frac{du_1}{U_1} = \frac{du_2}{U_2} = \cdots = \frac{du_{n+1}}{U_{n+1}};
\]
and if it happens that we have e.g.\ $U_1$ identically equal to zero, then the system becomes
\[
\frac{du_2}{U_2} = \frac{du_3}{U_3} = \cdots = \frac{du_{n+1}}{U_{n+1}},
\]
so that we have an integral $u_1 = c$, and then in the remaining equations substituting this value, or treating it, as constant, the system is reduced to one of $n-1$ equations.

Or again, if it happen that we have in the transformed system $m$ equations ($m < n$), say
\[
\frac{du_1}{U_1} = \frac{du_2}{U_2} = \cdots = \frac{du_{m+1}}{U_{m+1}},
\]
which are such that $U_1$, $U_2$, \dots, $U_{m+1}$ are functions of only the $m+1$ variables $u_1$, $u_2$, \dots, $u_{m+1}$, then the integration of the proposed system of $n$ equations depends on the integration in the first instance of a system of $m$ equations. It is to be observed that if the system of $m$ equations can be integrated, then the completion of the integration of the original system depends on the integration of a system of $n - m$ equations, and in this sense the original system of $n$ equations may be said to be broken up into two systems of $m$ equations and $n-m$ equations respectively: but \textit{non constat} that the system of $m$ equations admits of integration; and it is therefore more correct to say that, from the original system of the $n$ equations, there has been separated off a system of $m$ equations.

\textbf{113.} The bearing of the foregoing remarks on the problem of three bodies will presently appear. It will be seen that whereas the problem as it stood before Jacobi depends on a system of seven equations, it has been shown by him that there may be separated off from this a system of six equations.

\textbf{114.} Jacobi's memoir ``Sur l'\'elimination des Noeuds \&c.'' (1843).---The problem of the motion of three mutually attracting bodies is in the first instance reduced to that of the motion of two fictitious bodies (which may be considered as mutually attracting bodies, attracted by a fixed centre of force)$^{(1)}$.

\footnotetext[1]{This is the effect of Jacobi's reduction; but the explicit statement of the theorem, and actual replacement of the problem of the three bodies by that of the two bodies attracted to a fixed centre, is due to Bertrand (post, No.~117).}

In fact, in the original problem the centre of gravity of the three bodies moves uniformly in a right line, and it may without any real loss of generality be taken to be at rest; that is, if the $\xi$-coordinates of the three bodies are $\xi_1$, $\xi_2$, $\xi_3$, then $m_1\xi_1 + m_2\xi_2 + m_3\xi_3 = 0$, or $\xi_1$, $\xi_2$, $\xi_3$ may be taken to be linear functions of two quantities $x_1$ and $x_2$. And similarly for the $y$-coordinates and the $z$-coordinates respectively. And $(x_1, y_1, z_1)$, $(x_2, y_2, z_2)$ may be regarded as the coordinates of two bodies revolving about a fixed centre of force. Hence representing the differential coefficients in regard to the time by $x_1'$, \&c., and treating these as new variables, the equations of motion will assume the form
\[
\frac{dx_1}{X_1} = \frac{dy_1}{Y_1} = \frac{dz_1}{Z_1} = \frac{dx_2}{X_2} = \frac{dy_2}{Y_2} = \frac{dz_2}{Z_2} = \frac{dx_1'}{X_1'} = \frac{dy_1'}{Y_1'} = \frac{dz_1'}{Z_1'} = \frac{dx_2'}{X_2'} = \frac{dy_2'}{Y_2'} = \frac{dz_2'}{Z_2'},
\]
where $X_1$, $Y_1$, $Z_1$, $X_2$, $Y_2$, $Z_2$ are forces capable of representation by means of a force-function $U$. This is a system of twelve equations; but since $X_1$, $Y_1$, $Z_1$, $X_2$, $Y_2$, $Z_2$ are independent of the time, we may omit the equation $(=dt)$, and treat the system as one of eleven equations between the variables $x_1$, $y_1$, $z_1$, $x_2$, $y_2$, $z_2$, $x_1'$, $y_1'$, $z_1'$, $x_2'$, $y_2'$, $z_2'$: if this system were integrated, the determination of the time would then depend on a quadrature only. But for the system of eleven equations we have four integrals, viz., the integral of Vis Viva and the three integrals of areas, and the system is thus reducible to one of $(11 - 4 =)$ seven equations. This preliminary transformation in Jacobi's memoir explains the remark that the problem, as it stood before him, depended on a system of seven equations.

\textbf{115.} Jacobi remarks that in the transformed problem the three integrals of areas show (1) that the intersection of the planes of the orbits of the two bodies lie in a fixed plane, the invariable plane of the system; (2) that the inclinations of the planes of the two orbits to this fixed plane, and the parameters of the two orbits considered as variable ellipses, are four elements any two of which rigorously determine the two others.

And then choosing for variables the inclinations of the two orbits to the invariable plane, the two radius vectors, the angles which they form with the intersection of the planes of the two orbits, and lastly the angle between this intersection (being as already mentioned a line in the invariable plane) with a fixed line in the invariable plane, he finds that the last-mentioned angle entirely disappears from the system of differential equations, and is determined after their integration by a quadrature.

In this new form of the differential equations there is no trace of the nodes. The differential equations which determine the relative motion of the three bodies are reduced to five equations of the first order and one of the second order. The equations in question are the equations I.~to~VI. given at the end of the memoir.

It is to be remarked that $dt$ is not eliminated from these equations; the last of them is $\frac{d}{dt}(\mu r^2 + \mu_1 r_1^2) = 2U - 2h$; and if to reduce them to a system of equations of the first order we write $\frac{d}{dt}(\mu r^2 + \mu_1 r_1^2) = \theta$, and therefore $\frac{d\theta}{dt} = 2U - 2h$, the system may be presented in the form
\[
\frac{du}{U} = \frac{du_1}{U_1} = \frac{di}{I} = \frac{di_1}{I_1} = \frac{dr}{R} = \frac{dr_1}{R_1} = \frac{d\theta}{\Theta} = \frac{dt}{1},
\]
which if we do, and then omit the equation $(=dt)$, we have a system of six equations between the seven quantities $u$, $u_1$, $i$, $i_1$, $r$, $r_1$, $\theta$; when this is integrated, the equation $(=dt)$ gives the time by a quadrature; and finally, Jacobi's equation VII.\ ($d\Omega = \tan u \cdot dt$) gives by a quadrature the angle before referred to as disappearing from the system of equations I.~to~VI.

\textbf{116.} But when Jacobi says, ``Par suite on a fait cinq int\'egrations. Les int\'egrales connues n'\'etant qu'au nombre de quatre, on pourra donc dire que l'on a fait une int\'egration de plus dans le syst\`eme du monde. Je dis dans le syst\`eme du monde puisque la m\^eme m\'ethode s'applique \`a un nombre quelconque de corps,'' the language used is not, I think, quite accurate. It in fact appears from the memoir that it is only on the assumption of the integration of the system of six equations that, besides the integral of Vis Viva and the integrals of areas, the remaining two integrals are known; in fact after, but not before, the system of the order six has been integrated, the time $t$ and the angle $\Omega$ are each of them given by a quadrature.

\textbf{117.} Bertrand's ``M\'emoire sur l'int\'egration des \'Equations diff\'erentielles de la M\'ecanique'' (1852).---I have spoken of this memoir in No.~56 of my former Report. The course of investigation is the inquiry as to the integrals, which, combined according to Poisson's theorem with the integral of Vis Viva or any other given integral, give rise to an illusory result. But as regards the application made to the problem of three bodies, it will be more convenient to state from a different point of view the conclusions arrived at: and I may mention that when the author says ``Je parviens \dots\ \`a r\'eduire la question \`a l'int\'egration de six \'equations toutes du premier ordre, c'est-\`a-dire que j'effectue une int\'egration de plus que ne l'avait fait Jacobi,'' he seems to have overlooked that Jacobi's system of five equations of the first order and one of the second order really is, as above noticed, a system of the six equations with another equation which then gives the time by a quadrature, and that, at least as appears to me, he has not advanced the solution beyond the point to which it had been carried by Jacobi$^{(1)}$.

\footnotetext[1]{These remarks were communicated by me to M.~Bertrand---see my letter ``Sur l'int\'egration des \'Equations diff\'erentielles de la M\'ecanique,'' Comptes Rendus (1863)---and, in the answer he kindly sent me, he agrees that they are correct.}

\textbf{118.} Presenting Bertrand's results in the slightly different notation in which they are reproduced in Bour's memoir (post.\ No.~122), then if $(x, y, z)$, $(x_1, y_1, z_1)$ are the coordinates of the two bodies (the problem actually considered being, as with Jacobi, that of the motion of two bodies about a fixed centre of force), and representing the functions $x^2 + y^2 + z^2$, $x_1^2 + y_1^2 + z_1^2$, $m^2(x'^2 + y'^2 + z'^2)$, $m_1^2(x_1'^2 + y_1'^2 + z_1'^2)$, $m(xx' + yy' + zz')$, $m_1(x_1x_1' + y_1y_1' + z_1z_1')$, $m(xx_1 + yy_1 + zz_1)$, $(xx_1 + yy_1 + zz_1)^2$, $m^2m_1^2(x'x_1' + y'y_1' + z'z_1')$, $(xx_1 + yy_1 + zz_1)(x'x_1' + y'y_1' + z'z_1')$ by $u$, $u_1$, $v$, $v_1$, $w$, $w_1$, $r$, $r_1$, $q$, $s$ respectively, then the last-mentioned quantities are connected by a single geometrical relation, so that any one of them, say $s$, may be considered as a given function of the remaining nine. And the author in effect shows that the equations of motion give a system
\[
\frac{du}{U} = \frac{du_1}{U_1} = \frac{dv}{V} = \frac{dv_1}{V_1} = \frac{dw}{W} = \frac{dw_1}{W_1} = \frac{dr}{R} = \frac{dr_1}{R_1} = \frac{dq}{Q} = \frac{ds}{S},
\]
where $U$, $U_1$, \&c.\ are functions of the quantities $u$, $u_1$, $v$, \&c.\ Omitting from the system the equation $(=dt)$, there are eight equations between nine quantities; but there are two known integrals, viz., the integral of Vis Viva and the integral of principal moment (or sum of the squares of the integrals of areas); that is to say, the system is really a system of six equations.

\textbf{119.} Painvin, ``Recherche du dernier Multiplicateur \&c.'' (1854).---The author investigates the ultimate multiplier for two systems of differential equations:

$1^\circ$. The system of the equations I.~to~VI.\ in Jacobi's memoir ``Sur l'\'elimination des Noeuds \&c.'' (anti, No.~114). Writing in the equations $r_1 = r'$, $\frac{dr}{dt} = r''$, and treating $r'$, $r''$ as new variables, the system may be written in the form
\[
\frac{du}{U} = \frac{du_1}{U_1} = \frac{di}{I} = \frac{di_1}{I_1} = \frac{dr}{R} = \frac{dr_1}{R_1} = \frac{dr'}{R'} = \frac{dr''}{R''} = dt,
\]
which, omitting the equation $(=dt)$, is a system of seven equations between eight variables; and it is for this form of the system that the value of $M$ is determined, the result obtained being the simple and elegant one, $M = \frac{\sin i \sin i_1}{\theta}$. The system of seven equations has an integral which is in fact the equation V.\ of the system in Jacobi's form, so that it is really a system of six equations (anti, No.~115).

$2^\circ$. The system secondly discussed is Bertrand's system of nine equations (ante, No.~118). The multiplier $M$ is obtained under four different forms, $M = \frac{1}{\sqrt{BB' - AC}} = \frac{1}{\sqrt{AA'}} = \frac{1}{\sqrt{AZ + B}} = \frac{1}{mn}$ (I do not stop to explain the notation), the last of them being referred to as a result announced by M.~Bertrand in his course. But it is shown by M.~Bour in the memoir next referred to (post.\ No.~122), that the multiplier for the system in question can be obtained in a very much more simple manner, almost without calculation.

\textbf{120.} In connexion with Jacobi's theory of the elimination of the Nodes, I may refer to the investigations ``Application to the Problem of three Bodies,'' Nos.~84 to 96 of Donkin's memoir ``On a Class of Differential Equations \&c.'' Part II. The author remarks that his differential equations No.~93 afford an example of the so-called elimination of the Nodes, quite different however (in that they are merely transformations of the original differential equations of the problem without any integrations) from that effected by Jacobi.

\textbf{121.} It may be right to refer again in this place to the concluding part of \S~28 of Jacobi's memoir ``Nova Theoria Multiplicatoris \&c.'' (ante.\ No.~92), as bearing on the problem of three bodies.

\textbf{122.} Bour's ``M\'emoire sur le Probl\`eme des Trois Corps'' (1856).---The author remarks that Bertrand's system of equations have lost the remarkable form and the properties which characterize the ordinary equations for the solution of a dynamical problem. But by selecting eight new variables, functions of Bertrand's variables, the system may be brought back to the standard Hamiltonian form
\[
\frac{dq_i}{dt} = \frac{dH}{dp_i}, \qquad \frac{dp_i}{dt} = -\frac{dH}{dq_i},
\]
or to the form adopted by M.~Bour, of a partial differential equation
\[
\Sigma\left(\frac{dH}{dp_i}\frac{dC}{dq_i} - \frac{dH}{dq_i}\frac{dC}{dp_i}\right) = 0,
\]
and guiding himself by a theorem in relation to canonical integrals obtained in his memoir of 1855 (see No.~66 of my former Report), he finds by a somewhat intricate analysis the expressions of the eight new variables $p_1$, $p_2$, $p_3$, $p_4$, $q_1$, $q_2$, $q_3$, $q_4$. The results ultimately obtained are of a very remarkable and interesting form, viz.\ $H = \text{funct.}\,(p_1, p_2, p_3, p_4, q_1, q_2, q_3, q_4)$ is equal to the value it would have for motion in a plane, \textit{plus} a term admitting of a simple geometrical interpretation; and he thus arrives at the following theorem as a r\'esum\'e of the whole memoir, viz.,

``In order to integrate in the general case the problem of three bodies, it is sufficient to solve the case of motion in a plane, and then to take account of a disturbing function equal to the product of a constant depending on the areas by the sum of the moments of inertia of the bodies round a certain axis, divided by the square of the triangle formed by the three bodies.''

\textbf{123.} It may be remarked that the only given integral of the system of eight equations is the integral of Vis Viva, $H = \text{const.}$, and that using this equation to eliminate one of the variables, and omitting the equation $(=dt)$, we have, as in the solutions of Jacobi and Bertrand, a system of six equations between seven variables. As the equations are in the standard dynamical form, no investigation is needed of the multiplier $M$, which is given by Jacobi's general theory, and consequently when any five integrals of the six equations are given, the remaining integral can be obtained by a quadrature.

In the case of three bodies moving in a plane, the solution takes a very simple form, which is given in the concluding paragraph of the memoir.


\section*{Transformation of Coordinates}
\noindent\textsc{Article Nos.~124 to 141.}

\medskip

\textbf{124.} It may be convenient to remark at once that two sets of rectangular coordinates may be related to each other properly or improperly, viz., the axes to which they belong (considered as drawn from the origin in the positive directions) may be either capable, or else incapable, of being brought into coincidence. The latter relation, although of equal generality with the former one, may for the most part be disregarded; for by merely reversing the directions of the one set of axes, the improper is converted into the proper relation.

\textbf{125.} In the memoir ``Problema Algebraicum \&c.'' (1770) Euler proposes to himself the question ``Invenire novem numeros ita in quadratum disponendos ut satisfiat duodecem sequentibus conditionibus \&c.'', viz., substituting for $A$, $B$, $C$, \&c.\ the ordinary letters
\[
\begin{array}{ccc}
\alpha, & \beta, & \gamma, \\
\alpha', & \beta', & \gamma', \\
\alpha'', & \beta'', & \gamma'',
\end{array}
\]
the twelve conditions are
\begin{align*}
\alpha^2 + \alpha'^2 + \alpha''^2 &= 1, & \beta^2 + \beta'^2 + \beta''^2 &= 1, & \gamma^2 + \gamma'^2 + \gamma''^2 &= 1, \\
\alpha'^2 + \beta'^2 + \gamma'^2 &= 1, & \alpha''^2 + \beta''^2 + \gamma''^2 &= 1, & \alpha^2 + \beta^2 + \gamma^2 = 1, \\
\alpha\alpha' + \beta\beta' + \gamma\gamma' &= 0, & \alpha'\alpha'' + \beta'\beta'' + \gamma'\gamma'' &= 0, & \alpha\alpha'' + \beta\beta'' + \gamma\gamma'' = 0, \\
\alpha\alpha'' + \beta'\alpha'' + \gamma'\gamma'' &= 0, & \alpha\alpha' + \beta''\beta' + \gamma''\gamma' &= 0, & \alpha\beta + \alpha'\beta' + \alpha''\beta'' = 0.
\end{align*}
And he remarks that this is in fact the problem of the transformation of coordinates, viz., if we have
\[
X = \alpha x + \beta y + \gamma z, \quad Y = \alpha' x + \beta' y + \gamma' z, \quad Z = \alpha'' x + \beta'' y + \gamma'' z,
\]
then the first equations are such as to give identically
\[
X^2 + Y^2 + Z^2 = x^2 + y^2 + z^2.
\]

\textbf{126.} Assuming the first six equations, he shows by a direct analytical process that $0 = (\beta'\gamma'' - \beta''\gamma')^2$, or $\alpha = \pm(\beta'\gamma'' - \beta''\gamma')$; and taking the positive sign (for, as the numbers may be taken as well positively as negatively, there is nothing lost by doing so) $\alpha = \beta'\gamma'' - \beta''\gamma'$, which gives the system
\begin{align*}
\alpha &= \beta'\gamma'' - \beta''\gamma', & \beta &= \gamma'\alpha'' - \gamma''\alpha', & \gamma &= \alpha'\beta'' - \alpha''\beta', \\
\alpha' &= \beta''\gamma - \beta\gamma'', & \beta' &= \gamma''\alpha - \gamma\alpha'', & \gamma' &= \alpha''\beta - \alpha\beta'', \\
\alpha'' &= \beta\gamma' - \beta'\gamma, & \beta'' &= \gamma\alpha' - \gamma'\alpha, & \gamma'' &= \alpha\beta' - \alpha'\beta,
\end{align*}
and from these he deduces the second system of six equations.

The inverse system of equations
\[
X = \alpha x + \alpha' y + \alpha'' z, \quad Y = \beta x + \beta' y + \beta'' z, \quad Z = \gamma x + \gamma' y + \gamma'' z
\]
is not explicitly referred to.

\textbf{127.} He then satisfies the equations by means of trigonometrical substitutions, viz., assuming $\alpha = \cos\xi$, then $\alpha'^2 + \alpha''^2 = \sin^2\xi$, which is satisfied by $\alpha' = \sin\xi\cos\eta$, $\alpha'' = \sin\xi\sin\eta$, \&c., and he thus obtains for the coefficients a set of values involving the angles $\xi$, $\eta$, $\theta$, which are the same as those mentioned post, No.~130. And he shows how these formul\ae\ may be obtained geometrically by three successive transformations of two coordinates only. The remainder of the memoir relates to the analogous problem of the transformation of four or more coordinates.

\textbf{128.} I have analysed so much of Euler's memoir in order to show that it contains nearly the whole of the ordinary theory of the transformation of coordinates; the only addition required is the equation
\[
\begin{vmatrix}
\alpha & \beta & \gamma \\
\alpha' & \beta' & \gamma' \\
\alpha'' & \beta'' & \gamma''
\end{vmatrix} = \pm 1,
\]
where the sign $+$ gives $\alpha = \beta'\gamma'' - \beta''\gamma'$, \&c.\ (ut supra), but the sign $-$ would give $\alpha = -(\beta'\gamma'' - \beta''\gamma')$, \&c.

\textbf{129.} The distinction of the ambiguous sign is in fact the above-mentioned one of the proper and improper transformations; viz., for the sign $+$ the two sets of axes can, for the sign $-$ they cannot, be brought into coincidence: this very important remark was, I believe, first made by Jacobi in one of his early memoirs in Crelle's Journal, but I have lost the reference. As already mentioned, it is allowable to attend only to the proper transformation, and to consider the value of the determinant as being $= +1$; and this is in fact almost always done.

\textbf{130.} Euler's formul\ae\ involving the three angles are those which are ordinarily made use of in the problem of rotation and the problems of physical astronomy generally. It is convenient to take them as in the figure, viz.\ $\theta$, the longitude of node, $\phi$, the inclination, $\tau$, the angular distance of $X$ from node, and the formul\ae\ of transformation then are
\[
\begin{array}{c|ccc}
& X & Y & Z \\ \hline
x & \cos\tau\cos\theta - \sin\tau\sin\theta\cos\phi & \cos\tau\sin\theta + \sin\tau\cos\theta\cos\phi & \sin\tau\sin\phi \\
y & -\sin\tau\cos\theta - \cos\tau\sin\theta\cos\phi & -\sin\tau\sin\theta + \cos\tau\cos\theta\cos\phi & \cos\tau\sin\phi \\
z & \sin\theta\sin\phi & -\cos\theta\sin\phi & \cos\phi
\end{array}
\]

The foregoing very convenient algorithm, viz.\ the employment of
\[
\begin{array}{c|ccc}
& X & Y & Z \\ \hline
x & \alpha & \beta & \gamma \\
y & \alpha' & \beta' & \gamma' \\
z & \alpha'' & \beta'' & \gamma''
\end{array}
\]
to denote the system of equations
\[
X = \alpha x + \beta y + \gamma z, \quad \&\text{c.},
\]
is due to M.~Lam\'e in the memoir ``Du Mouvement de Rotation \&c.,'' \textit{Mem.\ de Berlin} for 1758 (pr.~1765).

\textbf{131.} But previously to the foregoing investigations, viz., Euler had obtained incidentally a very elegant solution of the problem of the transformation of coordinates; this is in fact identical with the next mentioned one, the letters $l$, $m$, $n$; $\lambda$, $\mu$, $\nu$ being used in the place of $\xi$, $\xi'$, $\xi''$; $\eta$, $\eta'$, $\eta''$.

\textbf{132.} In the memoir ``Formulae generales pro translatione \&c.'' (1775), Euler gives the following formul\ae\ for the transformation of coordinates, viz.\ if the position of the set of axes $XYZ$ in reference to the set $xyz$ is determined by
\[
\widehat{xX}, \widehat{yX}, \widehat{zX} = 90^\circ - \xi, \quad 90^\circ - \xi', \quad 90^\circ - \xi'',
\]
\[
\widehat{yXx}, \widehat{yXy}, \widehat{yXz} = \eta, \quad \eta', \quad \eta'',
\]
then the formul\ae\ of transformation are
\[
\begin{array}{c|ccc}
& X & Y & Z \\ \hline
x & \sin\xi & \cos\xi\sin\eta & \cos\xi\cos\eta \\
y & \sin\xi' & \cos\xi'\sin\eta' & \cos\xi'\cos\eta' \\
z & \sin\xi'' & \cos\xi''\sin\eta'' & \cos\xi''\cos\eta''
\end{array}
\]
with the following equations connecting the six angles, viz.\ if
\[
-A^2 = \cos(\eta' - \eta'')\cos(\eta'' - \eta)\cos(\eta - \eta'),
\]
then
\[
\tan\xi = \frac{-A}{\cos(\eta'' - \eta)}, \quad
\tan\xi' = \frac{-A}{\cos(\eta - \eta')}, \quad
\tan\xi'' = \frac{-A}{\cos(\eta' - \eta'')}.
\]

\textbf{133.} It is right to notice that these values of $\xi$, $\xi'$, $\xi''$ give the twelve equations $\alpha^2 + \beta^2 + \gamma^2 = 1$, \&c., but they do not give definitely $\alpha = \beta'\gamma'' - \beta''\gamma'$, \&c., but only $\alpha = \pm(\beta'\gamma'' - \beta''\gamma')$; that is, in the formul\ae\ in question the two sets of axes are not of necessity displacements the one of the other. In the same memoir Euler considers two sets of rectangular axes, and assuming that they are displacements the one of the other (this assumption is not made as explicitly as it should have been), he remarks that the one set may be made to coincide with the other set by means of a finite rotation about a certain axis (which may conveniently be termed the Resultant Axis). This consideration leads him to an equation which ought to $= 0$, be satisfied by the coefficients of transformation, but which he is not able to verify by means of the foregoing expressions in terms of $\xi$, $\xi'$, $\xi''$, $\eta$, $\eta'$, $\eta''$.

\textbf{134.} I remark that Euler's equation in fact is
\[
\begin{vmatrix}
\alpha - 1 & \beta & \gamma \\
\alpha' & \beta' - 1 & \gamma' \\
\alpha'' & \beta'' & \gamma'' - 1
\end{vmatrix} = 0,
\]
or, as it may be written,
\[
-(\beta'\gamma'' - \beta''\gamma') - (\gamma''\alpha - \gamma'\alpha'') - (\alpha\beta' - \alpha'\beta) + \alpha + \beta' + \gamma'' - 1 = 0,
\]
in which form it is an immediate consequence of the equations
\[
\begin{vmatrix}
\alpha & \beta & \gamma \\
\alpha' & \beta' & \gamma' \\
\alpha'' & \beta'' & \gamma''
\end{vmatrix} = 1, \quad \alpha = \beta'\gamma'' - \beta''\gamma', \quad \&\text{c.},
\]
which are true for the proper, but not for the improper transformation.

\textbf{135.} In the undated addition to the memoir, Euler states the theorem of the resultant axis as follows:---``Theorema, Quomodocunque sph\ae ra circa centrum suum convertatur, semper assignari potest diameter cujus directio in situ translato conveniat cum situ originali;'' and he again endeavours to obtain a verification of the foregoing analytical theorem.

\textbf{136.} The theory of the Resultant Axis was further developed by Euler in the memoir ``Nova Methodus Motum \&c.'' (1775), and by Lexell in the memoir ``Nonnulla theoremata generalia \&c.'' (1775): the geometrical investigations are given more completely and in greater detail in Lexell's memoir. The result is contained in the following system of formul\ae\ for the transformation of coordinates, viz., if $\alpha$, $\beta$, $\gamma$ are the inclinations of the resultant axis to the original set, and if $\phi$ is the rotation about the resultant axis, or say the resultant rotation, then we have
\[
\begin{array}{c|ccc}
& X & Y & Z \\ \hline
x & \cos^2\alpha + \sin^2\alpha\cos\phi & \cos\beta\cos\alpha(1 - \cos\phi) - \cos\gamma\sin\phi & \cos\gamma\cos\alpha(1 - \cos\phi) + \cos\beta\sin\phi \\
y & \cos\alpha\cos\beta(1 - \cos\phi) + \cos\gamma\sin\phi & \cos^2\beta + \sin^2\beta\cos\phi & \cos\gamma\cos\beta(1 - \cos\phi) - \cos\alpha\sin\phi \\
z & \cos\alpha\cos\gamma(1 - \cos\phi) - \cos\beta\sin\phi & \cos\beta\cos\gamma(1 - \cos\phi) + \cos\alpha\sin\phi & \cos^2\gamma + \sin^2\gamma\cos\phi
\end{array}
\]
Euler attempts, but not very successfully, to apply the formul\ae\ to the dynamical problem of the rotation of a solid body: he does not introduce them into the differential equations, but only into the integral ones, and his results are complicated and inelegant. The further simplification effected by Rodrigues was in fact required.

\textbf{137.} Jacobi's paper, ``Euleri formul\ae \&c.'' (1827), merely cites the last-mentioned result.

\textbf{138.} I find it stated in Lacroix's \textit{Differential Calculus}, t.~I.~p.~533, that the following system for the transformation of coordinates was obtained by Monge (no reference is given in Lacroix), viz., the system being as above,
\[
\begin{array}{ccc}
\alpha, & \beta, & \gamma, \\
\alpha', & \beta', & \gamma', \\
\alpha'', & \beta'', & \gamma'',
\end{array}
\]
and the quantities $\alpha$, $\beta'$, $\gamma''$ being arbitrary, then putting
\begin{align*}
1 + \alpha + \beta' + \gamma'' &= M, \\
1 + \alpha - \beta' - \gamma'' &= N, \\
1 - \alpha + \beta' - \gamma'' &= P, \\
1 - \alpha - \beta' + \gamma'' &= Q,
\end{align*}
so that $M + N + P + Q = 4$, we have
\begin{align*}
2\beta &= \sqrt{NP} + \sqrt{MQ}, & 2\gamma &= \sqrt{MP} + \sqrt{NQ}, & 2\alpha'' &= \sqrt{QN} + \sqrt{MP}, \\
2\beta' &= \sqrt{NP} - \sqrt{MQ}, & 2\gamma' &= \sqrt{PQ} - \sqrt{MN}, & 2\alpha' &= \sqrt{QN} - \sqrt{MP}.
\end{align*}
These are formul\ae\ very closely connected with those of Rodrigues.

\textbf{139.} The theory was perfected by Rodrigues in the valuable memoir ``Des lois g\'eom\'etriques \&c.'' (1840). Using for greater convenience $\lambda$, $\mu$, $\nu$ in the place of his $\frac{1}{2}\rho m$, $\frac{1}{2}\rho n$, $\frac{1}{2}\rho p$, he in effect writes
\[
\tan\tfrac{1}{2}\phi\,\cos\alpha = \lambda, \quad \tan\tfrac{1}{2}\phi\,\cos\beta = \mu, \quad \tan\tfrac{1}{2}\phi\,\cos\gamma = \nu,
\]
and this being so, the coefficients of transformation are
\begin{align*}
\frac{1 + \lambda^2 - \mu^2 - \nu^2}{1 + \lambda^2 + \mu^2 + \nu^2}, &\quad \frac{2(\nu + \mu\lambda)}{1 + \lambda^2 + \mu^2 + \nu^2}, \quad \frac{2(\mu\nu - \lambda)}{1 + \lambda^2 + \mu^2 + \nu^2}, \\
\frac{2(\lambda\mu - \nu)}{1 + \lambda^2 + \mu^2 + \nu^2}, &\quad \frac{1 - \lambda^2 + \mu^2 - \nu^2}{1 + \lambda^2 + \mu^2 + \nu^2}, \quad \frac{2(\lambda\nu + \mu)}{1 + \lambda^2 + \mu^2 + \nu^2}, \\
\frac{2(\lambda\nu + \mu)}{1 + \lambda^2 + \mu^2 + \nu^2}, &\quad \frac{2(\mu\nu - \lambda)}{1 + \lambda^2 + \mu^2 + \nu^2}, \quad \frac{1 - \lambda^2 - \mu^2 + \nu^2}{1 + \lambda^2 + \mu^2 + \nu^2},
\end{align*}
all divided by the common denominator $1 + \lambda^2 + \mu^2 + \nu^2$.

Conversely, if the coefficients of transformation are as usual represented by
\[
\begin{array}{ccc}
\alpha, & \beta, & \gamma, \\
\alpha', & \beta', & \gamma', \\
\alpha'', & \beta'', & \gamma'',
\end{array}
\]
then $\lambda$, $\mu$, $\nu$; $\lambda$, $\mu$, $\nu$ are respectively equal to
\[
\frac{\gamma - \beta''}{1 + \alpha + \beta' + \gamma''}, \quad
\frac{\alpha'' - \gamma}{1 + \alpha + \beta' + \gamma''}, \quad
\frac{\beta - \alpha'}{1 + \alpha + \beta' + \gamma''},
\]
each of them divided by $1 + \alpha + \beta' + \gamma''$.

The memoir contains very elegant formul\ae\ for the composition of finite rotations, and it will be again referred to in speaking of the kinematics of a solid body.

\textbf{140.} Sir W.~R.~Hamilton's first papers on the theory of quaternions were published in the years 1843 and 1844: the fundamental idea consists in the employment of the imaginaries $i$, $j$, $k$, which are such that
\[
i^2 = j^2 = k^2 = -1, \quad jk = -kj = i, \quad ki = -ik = j, \quad ij = -ji = k,
\]
whence also
\begin{align*}
&(w + ix + jy + kz)(w' + ix' + jy' + kz') \\
&\quad = ww' - xx' - yy' - zz' \\
&\qquad + i(wx' + w'x + yz' - y'z) \\
&\qquad + j(wy' + w'y + zx' - z'x) \\
&\qquad + k(wz' + w'z + xy' - x'y);
\end{align*}
so that representing the right-hand side by $W + iX + jY + kZ$, we have identically
\[
W^2 + X^2 + Y^2 + Z^2 = (w^2 + x^2 + y^2 + z^2)(w'^2 + x'^2 + y'^2 + z'^2).
\]
It is hardly necessary to remark that Sir W.~R.~Hamilton in his various publications on the subject, and in the \textit{Lectures on Quaternions}, Dublin, 1853, has developed the theory in detail, and has made the most interesting applications of it to geometrical and dynamical questions; and although the first explicit application of it to the present question may have been made in my own paper next referred to, it seems clear that the whole theory was in its original conception intimately connected with the notion of rotation.

\textbf{141.} Cayley, ``On certain Results relating to Quaternions'' (1845).---It is shown that Rodrigues' transformation formula may be expressed in a very simple manner by means of quaternions; viz., we have
\[
ix + jy + kz = (1 + i\lambda + j\mu + k\nu)^{-\frac{1}{2}}(iX + jY + kZ)(1 + i\lambda + j\mu + k\nu)^{\frac{1}{2}},
\]
where developing the function on the right-hand side, and equating the coefficients of $i$, $j$, $k$, we obtain the formul\ae\ in question. A subsequent paper, Cayley, ``On the application of Quaternions to the Theory of Rotation'' (1848), relates to the composition of rotations.


\section*{Principal Axes, and Moments of Inertia}
\noindent\textsc{Article Nos.~142---163.}

\medskip

\textbf{142.} The theorem of principal axes consists herein, that at any point of a solid body there exists a system of axes $Ox$, $Oy$, $Oz$, such that
\[
\int yz\,dm = 0, \quad \int zx\,dm = 0, \quad \int xy\,dm = 0.
\]
But this, the original form of the theorem, is a mere deduction from a general theory of the representation of the integrals
\[
\int x^2\,dm, \quad \int y^2\,dm, \quad \int z^2\,dm, \quad \int yz\,dm, \quad \int zx\,dm, \quad \int xy\,dm
\]
for any axes through the given origin by means of an ellipsoid depending on the values of these integrals corresponding to a given set of rectangular axes through the same origin.

\textbf{143.} If, for convenience, we write as follows, $M = \int dm$ the mass of the body, and
\[
A' = \int x^2\,dm, \quad B' = \int y^2\,dm, \quad C' = \int z^2\,dm, \quad F' = \int yz\,dm, \quad G' = \int zx\,dm, \quad H' = \int xy\,dm,
\]
and moreover
\begin{align*}
A &= \int(y^2 + z^2)\,dm, & B &= \int(z^2 + x^2)\,dm, & C &= \int(x^2 + y^2)\,dm, \\
F &= -\int yz\,dm, & G &= -\int zx\,dm, & H &= -\int xy\,dm,
\end{align*}
so that
\[
A = B' + C', \quad B = C' + A', \quad C = A' + B', \quad F = -F', \quad G = -G', \quad H = -H',
\]
then the ellipsoid which in the first instance presents itself for this purpose, and which Prof.~Price has termed the Ellipsoid of Principal Axes, but which I would rather term the ``Comomental Ellipsoid,'' is the ellipsoid
\[
(A', B', C', F', G', H'\,\widehat{\,}\,x, y, z)^2 = Mk^2,
\]
where $k$ is arbitrary, so that the absolute magnitude is not determined. But it is more usual, and in some respects better to consider in place thereof the ``Momental Ellipsoid'' (Cauchy, ``Sur les Moments d'Inertie,'' \textit{Exercices de Math\'ematique}, t.~II.\ pp.~93---103, 1827),
\[
(A, B, C, F, G, H\,\widehat{\,}\,x, y, z)^2 = Mk^2,
\]
or as it may also be written,
\[
(A' + B' + C')(x^2 + y^2 + z^2) - (A', B', C', F', G', H'\,\widehat{\,}\,x, y, z)^2 = Mk^2,
\]
which shows that the two ellipsoids have their axes, and also their circular sections, coincident in direction.

\textbf{144.} And there is besides this a third ellipsoid, the ``Ellipsoid of Gyration,'' which is the reciprocal of the momental ellipsoid in regard to the concentric sphere, radius $k$. The last-mentioned ellipsoid is given in magnitude, viz.\ if the body is referred to its principal axes, then putting $A = Ma^2$, $B = Mb^2$, $C = Mc^2$, the equation of the ellipsoid of gyration is
\[
\frac{x^2}{a^2} + \frac{y^2}{b^2} + \frac{z^2}{c^2} = k^2.
\]
The axes of any one of the foregoing ellipsoids coincide in direction with the principal axes of the body, and the magnitudes of the axes lead very simply to the values of the principal moments $A$, $B$, $C$.

\textbf{145.} The origin has so far been left arbitrary: in the dynamical applications, this origin is in the case of a solid body rotating about a fixed point, the fixed point; and in the case of a free body, the centre of gravity. But the values of the coefficients $(A, B, C, F, G, H)$, or $(A', B', C', F', G', H')$, corresponding to any given origin whatever, are very easily expressed in terms of the coordinates of this origin, and the values of the corresponding coefficients for the centre of gravity as origin; or, what is the same thing, any one of the ellipsoids for the given origin may be geometrically constructed by means of the ellipsoid for the centre of gravity. The geometrical theory, as regards the magnitudes of the axes, does not appear to have been anywhere explicitly enunciated; as regards their direction, it is comprised in the theorem that the directions at any point are the three rectangular directions at that point in regard to the ellipsoid of gyration for the centre of gravity$^{(1)}$, post, No.~159. The notion of the ellipsoids, and of the relation between the ellipsoids at a given point and those at the centre of gravity, once established, the theory of principal axes and moments of inertia becomes a purely geometrical one.

\footnotetext[1]{The rectangular directions at a point in regard to an ellipsoid are the directions of the axes of the circumscribed cone, or, what is the same thing, they are the directions of the normals to the three quadric surfaces, confocal with the given ellipsoid, which pass through the given point. The theory of confocal surfaces appears to have been first given by Chasles, Note XXXI.\ of the \textit{Aper\c{c}u Historique} (1837).}

\textbf{146.} The existence of principal axes was first established by Segner in the work \textit{Specimen Theori\ae\ Turbinum}, Halle (1755), where, however, it is remarked that Euler had said something on the subject in the (Berlin) \textit{Memoirs} for 1749 and 1750 (post.\ No.\ 167), and had constructed a new mechanical principle, but without pursuing the question. Segner's course of investigation is in principle the same as that now made use of, viz.\ a principal axis is defined to be an axis such that when a body revolves round it the forces arising from the rotation have no tendency to alter the position of the axes. It is first shown that there are systems of axes $x$, $y$, $z$ such that $\int xz\,dm = 0$, and then, in reference to such a set of axes, the position of a principal axis, say the axis of $X$, is determined by the conditions
\[
\int XY\,dm = 0, \quad \int XZ\,dm = 0,
\]
being taken to be $t = \frac{X}{x}$, $\tau = \frac{Y}{y}$ (where $\alpha$, $\beta$, $\gamma$ being the unknown quantities, the inclinations of the principal axis to those of $x$, $y$, $z$), and then putting $A' = \int x^2\,dm$, \&c.\ ($F' = 0$ by hypothesis), Segner's equations for the determination of $t$, $\tau$ are
\begin{align*}
G't + (C' - A')\tau - G' &= 0, \\
H't + (C' - B')t\tau - G'\tau^2 + H'\tau &= 0,
\end{align*}
the second of which gives
\[
\tau = \frac{H't}{G'H'' - G'(A' - B')t + \{(B' - C')(C' - A') - G''^2 - H''^2\}t + G'(B' - C'')},
\]
which being a cubic equation shows that there are three principal axes; and it is afterwards proved that these are at right angles to each other.

\textbf{147.} To show the equivalence of Segner's solution to the modern one, I remark that if $M = \int X^2\,dm$, we have
\begin{align*}
(A' - u)t + H'\tau + G' &= 0, \\
H't + (B' - u)\tau + F' &= 0, \\
G't + F'\tau + (C' - u) &= 0,
\end{align*}
whence
\begin{align*}
t^2 : \tau^2 : 1 : t : \tau : t\tau &= \\
&\hspace{-2cm} B'C' - F'^2 - (B' + C')u + u^2 : C'A' - G'^2 - (C' + A')u + u^2 : A'B' - H'^2 - (A' + B')u + u^2 \\
&\hspace{-2cm} : G'F' - A'F' + F'u : H'F' - B'G' + G'u : F'G' - C'H' + H'u.
\end{align*}
Or putting therein $F' = 0$,
\begin{align*}
t^2 : \tau^2 : 1 : t : \tau : t\tau &= \\
&\hspace{-2cm} B'C' - (B' + C')u + u^2 : C'A' - G'^2 - (C' + A')u + u^2 : A'B' - H'^2 - (A' + B')u + u^2 \\
&\hspace{-2cm} : -FG' : +H'G'u : G'H' + (-C'H' + H'u)u,
\end{align*}
by means of which Segner's equations may be verified. I have given this analysis, as the first solution of such a problem is a matter of interest.

\textbf{148.} There is little if anything added to Segner's results by the memoir, Euler, ``Recherches sur la Connaissance M\'ecanique des Corps'' (1758), which is introductory to the immediately following one on Rotation.

\textbf{149.} Relating to the theory of principal axes we have Binet's ``M\'emoire sur les axes conjugu\'es \&c.,'' (1813). The author proposes to make known the new systems of axes which he calls conjugate axes, which, when they are at right angles to each other, coincide with the principal axes; viz.\ considering the sum of the molecules each into its distance from a plane, such distance being measured in the direction of a line, then (the direction of the line being given) of all the planes which pass through a given point, there is one for which the sum in question is a minimum, and this plane is said to be conjugate to the given line, and from the notion of a line and conjugate plane he passes to that of a system of conjugate axes. The investigation (which is throughout an elegant one) is conducted analytically; the coordinates made use of are oblique ones, and the formul\ae\ are thus rendered more complicated than they would have been; in referring to them it will be convenient to make the axes rectangular.

\textbf{150.} One of the results is the well-known equation
\[
(A' - \theta)(B' - \theta)(C' - \theta) - F'^2(A' - \theta) - G'^2(B' - \theta) - H'^2(C' - \theta) + 2F'G'H' = 0;
\]
which, if $x_1$, $y_1$, $z_1$ are the principal axes, has for its roots $\int x_1^2\,dm$, $\int y_1^2\,dm$, $\int z_1^2\,dm$.

And the equations (1), p.~49, taking therein the original axes as rectangular, are
\begin{align*}
[\mathfrak{A}' - \mathfrak{K}']\cos\alpha + \mathfrak{H}'\cos\beta + \mathfrak{G}'\cos\gamma &= 0, \\
\mathfrak{H}'\cos\alpha + [\mathfrak{B}' - \mathfrak{K}']\cos\beta + \mathfrak{F}'\cos\gamma &= 0, \\
\mathfrak{G}'\cos\alpha + \mathfrak{F}'\cos\beta + [\mathfrak{C}' - \mathfrak{K}']\cos\gamma &= 0,
\end{align*}
where $\mathfrak{A}'$, $\mathfrak{B}'$, $\mathfrak{F}'$, $\mathfrak{G}'$, $\mathfrak{H}'$, \&c.\ denote the reciprocal coefficients, $\mathfrak{A}' = B'C' - F'^2$, \&c., and $\mathfrak{K}'$ is the discriminant $= A'B'C' - A'F'^2 - B'G'^2 - C'H'^2 + 2F'G'H'$: this is a symmetrical system of equations for finding $\cos\alpha : \cos\beta : \cos\gamma$, less simple however than the modern form (post, No.~154), the identity of which with Binet's may be shown without difficulty.

\textbf{151.} Another result (p.~57) is that if the original axes are principal axes, and if $Ox$, $Oy$, $Oz$ are the principal axes through a point the coordinates whereof are $f$, $g$, $h$, and if $O_x = ($say$)\int x_1^2\,dm$, then we have
\[
\frac{f^2}{O_x - A'} + \frac{g^2}{O_x - B'} + \frac{h^2}{O_x - C'} = \frac{1}{M},
\]
(in which I have restored the mass $M$, which is put equal to unity), so that if $O_x$ have a given constant value, the locus of the point is a quadric surface, the nature whereof will depend on the value of $O_x$. The surfaces in question are confocal with each other (and with the imaginary surface $\frac{f^2}{A'} + \frac{g^2}{B'} + \frac{h^2}{C'} = \frac{1}{M}$), and the general theorem in regard to them is, I think, due to Chasles: the special theorem for the ellipsoid of gyration was, however, given by Binet.

\textbf{152.} An interesting theorem of Binet's relates to the momental ellipsoid at any point of a principal plane: taking the fixed point in the plane of $xy$, then $h = 0$, and we have
\[
\frac{f^2}{O_x - A'} + \frac{g^2}{O_x - B'} = \frac{1}{M} - \frac{0}{O_x - C'} = \frac{1}{M},
\]
which may be written
\[
\frac{f^2}{\frac{1}{M}O_x - \frac{A'}{M}} + \frac{g^2}{\frac{1}{M}O_x - \frac{B'}{M}} = 1,
\]
and thus shows that the section by the plane of $xy$ of the system of quadric surfaces is a system of confocal conics; the foci thereof are the principal axes in the plane of $xy$, and the squares of the semiaxes thereof are $\frac{A'}{M}$ and $\frac{B'}{M}$, that is, the squares of the principal radii of gyration in the plane of $xy$.

\textbf{153.} Another theorem relates to the section of the ellipsoid of gyration by a principal plane. Taking this to be the plane of $xy$, and writing $A = Ma^2$, $B = Mb^2$, $C = Mc^2$, the equation of the ellipsoid of gyration is $\frac{x^2}{a^2} + \frac{y^2}{b^2} + \frac{z^2}{c^2} = k^2$; the section by the plane of $xy$ is $\frac{x^2}{a^2} + \frac{y^2}{b^2} = k^2$, an ellipse the squares of whose semiaxes are $a^2$, $b^2$, that is, the principal radii of gyration in the plane of $xy$. It may be added, that at any point of the plane of $xy$ the section of the ellipsoid of gyration is an ellipse having its foci at the principal axes (or centres of gyration) in that plane; and hence also the theorem that the difference of the squares of the radii of gyration in any two directions at right angles to each other in a principal plane is constant, being equal to the difference of the squares of the principal radii of gyration in that plane.

\textbf{154.} To complete the theory, I give the following proof of the theorem for the principal axes through any point of a solid body. Starting from the ellipsoid of gyration for the centre of gravity, the axes at the point in question are (ante No.~145) the rectangular directions at the point in regard to the ellipsoid. Now the rectangular directions at a point in regard to an ellipsoid coincide with the normals to the three confocal quadric surfaces through the point. Hence the directions of the principal axes at the point in question are those of the normals to the three confocal quadric surfaces through that point; and if we determine, as above, the axes of the confocal ellipsoid through the point, these will be the principal axes in question. It is to be observed that the axes of the confocal ellipsoid are not the same thing as the rectangular directions, and that the theorem does not apply to the axes of the confocal ellipsoid, but only to the normals thereof. The rectangular directions at any point in regard to an ellipsoid, and the normals to the confocal quadric surfaces through the point, are identical; and they coincide also with the axes of the cone circumscribed about the ellipsoid, having its vertex at the point. The notion of the rectangular directions is, I think, the most convenient one for the statement of the theorem.

\textbf{155.} The analytical expression of the theorem is as follows:---The equations of the confocal surfaces being
\[
\frac{x^2}{a^2 + \theta} + \frac{y^2}{b^2 + \theta} + \frac{z^2}{c^2 + \theta} = k^2,
\]
these determine for a given point $(f, g, h)$ three values $\theta_1$, $\theta_2$, $\theta_3$; and if $(l, m, n)$ are the direction cosines of an axis, then the equations for the determination of $l : m : n$ are
\begin{align*}
\left(\frac{f^2}{(a^2 + \theta_1)^2} + \frac{g^2}{(b^2 + \theta_1)^2} + \frac{h^2}{(c^2 + \theta_1)^2}\right) l &= \frac{f}{a^2 + \theta_1}, \\
\left(\frac{f^2}{(a^2 + \theta_1)^2} + \frac{g^2}{(b^2 + \theta_1)^2} + \frac{h^2}{(c^2 + \theta_1)^2}\right) m &= \frac{g}{b^2 + \theta_1}, \\
\left(\frac{f^2}{(a^2 + \theta_1)^2} + \frac{g^2}{(b^2 + \theta_1)^2} + \frac{h^2}{(c^2 + \theta_1)^2}\right) n &= \frac{h}{c^2 + \theta_1},
\end{align*}
for $\theta = \theta_1$, and two other similar systems for $\theta = \theta_2$ and $\theta = \theta_3$; these give the directions of the three principal axes at the point $(f, g, h)$.

\textbf{156.} I may remark that the foregoing investigations belong to a purely geometrical theory which I call the ``Th\'eorie des Directions Rectangulaires''; the theory might, I think, be advantageously introduced into the ordinary treatises on solid geometry.

\textbf{157.} Another theorem relating to the ellipsoid of gyration is the following:---If through any point $(f, g, h)$ we draw a line meeting the ellipsoid of gyration in the points $P$, $Q$, and if $R$ be the radius vector parallel to this line drawn from the centre to the ellipsoid, then the product $OP \cdot OQ$ (where $O$ is the point $(f, g, h)$) varies inversely as $R^2$; or, what is the same thing, the product $OP \cdot OQ \cdot R^2$ is constant for all directions through the point. This theorem is due to MacCullagh.

\textbf{158.} I do not stop to consider the various other theorems relating to moments of inertia and principal axes which are to be found in the works of Poinsot, Chasles, and other writers.

\textbf{159.} The theorem of the principal axes at any point, viz.\ that the directions thereof are the rectangular directions at the point in regard to the ellipsoid of gyration for the centre of gravity, has been already referred to, No.~145. It may be remarked that the theorem includes as a particular case the theorem for the principal axes through the centre of gravity itself; for at the centre of gravity the ellipsoid of gyration becomes the central ellipsoid of gyration, and the rectangular directions at the centre of gravity in regard to this ellipsoid are of course the axes thereof.

\textbf{160.} The ellipsoid of gyration at any point is connected with the central ellipsoid of gyration by a simple relation: taking the central ellipsoid to be $\frac{x^2}{a^2} + \frac{y^2}{b^2} + \frac{z^2}{c^2} = k^2$, then the ellipsoid of gyration at the point $(f, g, h)$ is at once expressible in terms of $f$, $g$, $h$, and the coefficients $a^2$, $b^2$, $c^2$; the equation may be written
\[
\frac{x^2}{a^2} + \frac{y^2}{b^2} + \frac{z^2}{c^2} = k^2 + \frac{f^2}{a^2} + \frac{g^2}{b^2} + \frac{h^2}{c^2} - \left(\frac{fx}{a^2} + \frac{gy}{b^2} + \frac{hz}{c^2}\right)^2 \bigg/ \left(\frac{f^2}{a^2} + \frac{g^2}{b^2} + \frac{h^2}{c^2}\right),
\]
and from this the axes of the ellipsoid, that is, the principal radii of gyration at the point, may be found. But as regards the directions only, these are, as has been seen, the rectangular directions at the point in regard to the central ellipsoid of gyration.

\textbf{161.} The investigation of the principal axes by means of the ellipsoid of gyration, although in some respects the most complete and elegant method, is not the one originally employed. The first direct analytical solution of the problem is that of Segner (1755), already referred to, No.~146; and it is worthy of remark that the method employed by Segner is in principle identical with that which is now in common use, viz.\ the problem is made to depend on the solution of a cubic equation.

\textbf{162.} The modern form of the solution is as follows:---If $A'$, $B'$, $C'$, $F'$, $G'$, $H'$ are the integrals $\int x^2\,dm$, \&c., referred to any system of rectangular axes through the point, then the determination of the principal axes depends on the solution of the cubic equation
\[
\begin{vmatrix}
A' - \theta & H' & G' \\
H' & B' - \theta & F' \\
G' & F' & C' - \theta
\end{vmatrix} = 0,
\]
or, expanded,
\[
(A' - \theta)(B' - \theta)(C' - \theta) - F'^2(A' - \theta) - G'^2(B' - \theta) - H'^2(C' - \theta) + 2F'G'H' = 0.
\]
If the roots are $\theta_1$, $\theta_2$, $\theta_3$, then the direction cosines $(l, m, n)$ of the corresponding principal axis are given by
\[
(A' - \theta_1)l + H'm + G'n = 0, \quad H'l + (B' - \theta_1)m + F'n = 0, \quad G'l + F'm + (C' - \theta_1)n = 0,
\]
with two other similar systems for $\theta_2$ and $\theta_3$.

\textbf{163.} The foregoing cubic equation has three real roots, and the corresponding axes are at right angles to each other. That the roots are real may be shown in various ways; the most direct proof is that the cubic is the discriminant of the quadratic form $(A', B', C', F', G', H'\,\widehat{\,}\,l, m, n)^2$, which being a positive definite form, its discriminant has three real roots. The theory of principal axes is thus made to depend on the theory of quadratic forms, and the whole investigation is conducted by purely algebraical methods.


\section*{Rotation of a Solid Body}
\noindent\textsc{Article Nos.~164 to 205.}

\medskip

\textbf{164.} The problem of the rotation of a solid body about a fixed point is one of the most interesting and important in dynamics. The first complete solution is due to Euler, who in the memoir ``Du Mouvement de Rotation des Corps Solides autour d'un Point Variable'' (1758) established the equations of motion and obtained the integrals of Vis Viva and of areas. But the complete solution, expressing the nine direction cosines in terms of the time, was not effected until the discoveries of Jacobi (1839) relating to elliptic functions.

\textbf{165.} The equations of motion, as given by Euler, are
\[
A\frac{dp}{dt} + (C - B)qr = L, \quad B\frac{dq}{dt} + (A - C)rp = M, \quad C\frac{dr}{dt} + (B - A)pq = N,
\]
where $A$, $B$, $C$ are the principal moments of inertia, $p$, $q$, $r$ the angular velocities about the principal axes, and $L$, $M$, $N$ the moments of the external forces about the same axes. In the case where there are no external forces, these become
\[
A\frac{dp}{dt} + (C - B)qr = 0, \quad B\frac{dq}{dt} + (A - C)rp = 0, \quad C\frac{dr}{dt} + (B - A)pq = 0.
\]

\textbf{166.} The two integrals in this case are those of Vis Viva and of areas:
\[
Ap^2 + Bq^2 + Cr^2 = 2h, \qquad A^2p^2 + B^2q^2 + C^2r^2 = k^2,
\]
where $h$ and $k$ are constants. From these equations, together with the third equation of motion, the values of $p$, $q$, $r$ may be expressed in terms of elliptic functions of the time. The actual expressions were first obtained by Jacobi in 1839.

\textbf{167.} Euler's investigations on the rotation of a solid body extend over a long series of years. The memoir of 1758, already referred to, contains the complete establishment of the equations of motion and the proof of the existence of the two integrals; but the further integration was not effected. In the memoir ``Problema Algebraicum propositum publice'' (1770), Euler considered the transformation of coordinates, which is intimately connected with the problem; and in the memoir ``Nova Methodus Motum Corporum rigidorum determinandi'' (1775), he returned to the problem of rotation, obtaining results which, although interesting, do not constitute a complete solution.

\textbf{168.} The first step towards the complete solution was made by Lagrange in the \textit{M\'ecanique Analytique} (1788), who showed that the equations of motion could be expressed in terms of the angles which determine the position of the body, and who obtained the integrals in a form suitable for further development. But the actual integration in terms of elliptic functions required the discoveries of Abel and Jacobi.

\textbf{169.} Poinsot's geometrical representation of the motion (1834) is one of the most beautiful results in the whole theory. The motion is represented by the rolling of the ellipsoid of gyration on a fixed plane, the distance of the plane from the centre being proportional to the angular velocity. This representation, although it does not give the analytical expressions for the motion, affords a complete and elegant geometrical description thereof.

\textbf{170.} The analytical solution was completed by Jacobi in the memoir ``Sur la Rotation d'un Corps'' (1839), and in the fundamental memoirs on elliptic functions in Crelle's Journal. The solution is expressed by means of the elliptic functions sn, cn, dn, and involves the introduction of the functions $\Theta$ and $H$, which Jacobi invented for this purpose. The expressions for the direction cosines are of the form
\[
\alpha = \frac{H(u + ia)}{\Theta(u)}, \quad \&\text{c.},
\]
and the time $t$ is proportional to the argument $u$.

\textbf{171.} Jacobi's solution was considerably simplified by the introduction of the notation of elliptic functions with the modulus $k$, and by the use of the addition theorems and other properties of these functions. The solution has been presented in various forms by different writers; among the most important are those of Rosenhain, Richelot, and others.

\textbf{172.} The case of two equal moments of inertia, say $A = B$, is of special interest, as it includes the case of a body of revolution. In this case the elliptic functions degenerate into circular functions, and the solution takes a much simpler form. The motion is then represented by the rolling of a spheroid on a fixed plane, and the precession and nutation of the axis of the body are easily determined.

\textbf{173.} The case of three equal moments of inertia, $A = B = C$, is the simplest of all; the body moves as if it were a sphere, and the axis of rotation remains fixed in direction in space. This is the case of the free rotation of a sphere or cube about its centre.

\textbf{174.} Another particular case of interest is that in which the angular momentum is the least or greatest possible for the given value of the Vis Viva; in this case the motion is uniform, and the axis of rotation is a principal axis. This is the case of steady rotation about a principal axis, which is stable for the axes of greatest and least moment, and unstable for the intermediate axis.

\textbf{175.} The stability of the rotation about a principal axis was first investigated by Euler, who showed that the rotation about the axes of greatest and least moment is stable, while that about the intermediate axis is unstable. The analytical proof depends on the consideration of small oscillations about the state of steady motion; if the axis be that of greatest or least moment, the oscillations remain small, but if it be the intermediate axis, they increase indefinitely.

\textbf{176.} The motion of a solid body subject to external forces is in general much more complicated than that of the free body. The most important case is that in which the body is acted on by gravity, and rotates about a fixed point not coinciding with the centre of gravity; this is the case of the heavy rigid body, or the ``top.'' The complete integration of this problem has not been effected, except in certain particular cases.

\textbf{177.} The particular cases in which the problem of the heavy top admits of complete solution are:---1.\ The case of Euler and Poinsot, where the fixed point coincides with the centre of gravity; 2.\ The case of Lagrange, where two of the principal moments are equal, and the centre of gravity lies on the axis of symmetry; 3.\ The case of Kowalevski, where two of the principal moments are equal to one-fourth of the third, and the centre of gravity lies in the plane perpendicular to the axis of the unequal moment.

\textbf{178.} The case of Lagrange (the symmetrical top) is the one of greatest physical interest, as it includes the motion of the spinning top and the precession and nutation of the earth. The solution is effected by means of elliptic functions, and the motion consists of a steady precession combined with a periodic nutation.

\textbf{179.} The case of Kowalevski (1888) was a remarkable discovery, showing that there exists a third case in which the problem admits of complete integration. The solution is more complicated than that of Lagrange, and involves hyperelliptic functions; but the existence of this case is of great theoretical importance.

\textbf{180.} The investigations of Kowalevski were continued by various writers, among whom may be mentioned Tchapliguine, Husson, and others. The problem of the complete integration of the equations of motion of a heavy rigid body remains, however, one of the most difficult in dynamics, and it is not known whether there are any other cases besides those already discovered.

\textbf{181.} The equations of motion of a heavy rigid body may be written in various forms, according to the variables employed. The most usual form is that of Euler, in terms of the angular velocities $p$, $q$, $r$ about the principal axes; but they may also be expressed in terms of the Eulerian angles $\theta$, $\phi$, $\psi$, or in terms of the parameters of Rodrigues, or by means of quaternions.

\textbf{182.} The use of quaternions in the theory of rotations, as developed by Hamilton and Tait, affords a very elegant method of treating the problem. The equation of motion may be written in the form
\[
\frac{d\rho}{dt} = \tfrac{1}{2}\rho\omega,
\]
where $\rho$ is the quaternion representing the position of the body, and $\omega$ is the vector quaternion representing the angular velocity. This form of the equation is particularly adapted to the treatment of the composition of rotations.

\textbf{183.} The theory of the composition of finite rotations, as developed by Rodrigues, Hamilton, and Cayley, is intimately connected with the problem of the rotation of a solid body. The composition of two finite rotations about different axes is equivalent to a single rotation about a certain axis; and the determination of the resultant rotation may be effected by various methods, among which the quaternion method is perhaps the most elegant.

\textbf{184.} The kinematical representation of the motion of a solid body, as given by Poinsot, has been already referred to, No.~169. Analytical representations of a similar character have been given by various writers; among the most important are the ``moving axes'' of Thomson and Tait, and the ``screws'' of Ball.

\textbf{185.} The method of moving axes, as developed by Thomson and Tait, affords a very powerful method of treating the kinematics and dynamics of a solid body. The axes are supposed to move with the body, and the components of velocity and angular velocity are referred to them; the equations of motion then take a form which is particularly adapted to the treatment of problems involving the rotation of a solid body.

\textbf{186.} Rueb's memoir (1834) on the rotation of a solid body is of great historical importance, as it contains the first successful attempt to express the direction cosines in terms of elliptic functions. The memoir is entitled ``De motu gyratorio corporis rigidi'' and was published in the \textit{Commentationes Societatis Regiae Scientiarum Gottingensis}. The solution, although incomplete in some respects, prepared the way for the complete solution of Jacobi.

\textbf{187.} In Rueb's solution, the direction cosines are expressed by means of certain integrals involving the square root of a quartic function; these integrals are in fact elliptic integrals of the first kind, but Rueb did not reduce them to the standard form. The reduction was effected by Jacobi, who thus completed the solution.

\textbf{188.} The connection between the problem of rotation and the theory of elliptic functions is one of the most interesting in the whole range of mathematics. The elliptic functions arise in the most natural manner from the attempt to express the direction cosines in terms of the time; and the theory of these functions receives its most beautiful applications in the solution of this problem.

\textbf{189.} The further development of the theory of elliptic functions, as given by Jacobi in the \textit{Fundamenta Nova} (1829), and in his later memoirs, was largely influenced by the requirements of the problem of rotation. The functions $\Theta$ and $H$, which are fundamental in the theory, were invented for the purpose of expressing the direction cosines in the most elegant form.

\textbf{190.} The solution of the problem of rotation has been presented in various forms by different writers. Among the most important of these forms are those of Jacobi (1839), Rosenhain (1846), Richelot (1850), and others. Each of these writers has contributed something to the simplification or extension of the solution.

\textbf{191.} Rosenhain's memoir (1846) contains a very elegant treatment of the problem, based on the theory of the hyperelliptic functions. The solution is expressed in terms of the quotients of the theta-functions, and the expressions are symmetrical and elegant. But the solution is less directly applicable to the physical problem than that of Jacobi.

\textbf{192.} Richelot's memoir (1850) contains an elaborate investigation of the problem, with especial reference to the case of unequal moments of inertia. The solution is expressed in terms of the elliptic functions, and the author has paid particular attention to the reduction of the integrals to the standard form. The memoir also contains various interesting particular cases.

\textbf{193.} The solution by Jacobi, \S~27 of the memoir ``Theoria Novi Multiplicatoris'' (1845), is given as an application of the general theory, the author remarking that, as well in this question as in the problem of the two fixed centres, he purposely employed a somewhat inartificial analysis, in order to show that the principle (of the Ultimate Multiplier) would lead to the result without any special artifices. The principal axes are referred to the axes fixed in space by the ordinary three angles (here called $q_1$, $q_2$, $q_3$), and the solution is carried so far as to give the integral equations, without any reduction of the integrals contained in them to elliptic integrals. The solution is, however, in so far remarkable that the integrations are effected without the aid of the invariable plane.

\textbf{194.} Cayley, ``On the Rotation of a Solid Body \&c.'' (1846).---It appeared desirable to obtain the solution by means of the quantities $\lambda$, $\mu$, $\nu$, without the assistance of the invariable plane, and Jacobi's discovery of the theorem of the Ultimate Multiplier induced me to resume the problem, and at least attempt to bring it so far as to obtain a differential equation of the first order between two variables only, the multiplier of which could be obtained theoretically by Jacobi's discovery. The choice of two new variables to which the equations of the problem led me, enabled me to effect this in a simple manner; and the differential equation which I finally obtained turned out to be integrable \textit{per se}, so that the laborious process of finding the multiplier became unnecessary.

\textbf{195.} The new variables $\Omega$, $v$ have the following geometrical significations: $\Omega = I\tan\frac{1}{2}\theta\cos J$, where $I$ is the principal moment ($A^2p^2 + B^2q^2 + C^2r^2 = I^2$), $\theta$ (as before) the angle of resultant rotation, and $J$ is the inclination of the resultant axis to the invariable axis; and $v = I^2\cos^2\frac{1}{2}\theta \cdot J$, where if we imagine a line $AQ$ having the same position relatively to the axes fixed in space that the invariable axis has to the principal axes of the body, then $J$ is the inclination of this line to the invariable axis. It is found that $p$, $q$, $r$ are functions of $v$ only, whereas $\lambda$, $\mu$, $\nu$ contain besides the variable $\Omega$. In obtaining these relations, there occurs a singular relation $\Omega^2 = kv - 1$, which may also be written $1 + \tan^2\frac{1}{2}\theta\cos^2 J = \text{const.}$, where the geometrical significations of the quantities $I$, $J$ have just been explained.

The final results are that the time $t$, and the arc $\tan^{-1}\frac{\Omega}{\sqrt{kv - 1}}$ are each of them expressible as the integrals of certain algebraical functions of $v$. There might be some interest in comparing the results with those of Euler's memoir of 1758, where the principal axes are also referred to an arbitrary system of axes fixed in space; but I was not then acquainted with Euler's memoir.

The concluding part of the paper relates to the determination of the variations of the constants in the disturbed problem.

\textbf{196.} Cayley, ``Note on the Rotation of a Solid of Revolution'' (1849), shows the simplification produced in the formul\ae\ of the last-mentioned memoir in the case where two of the moments of inertia are equal, say $A = B$.

\textbf{197.} Jacobi's final solution of the problem of Rotation was given without demonstration in the letter to the Academy of Sciences at Paris; the demonstration is added in the memoir, Crelle, t.~XXXIX.~(1849). The fundamental idea consists in taking in the invariable plane, instead of the fixed axes $xy$, a set of axes $xy$ revolving with uniform velocity, such that the angular distance of the axis of $x$ from its initial position is precisely $= -n'u$; and consequently if $\theta'$ be the longitude of the node of the equator on the invariable plane, measured from the moveable axis of $x$ as the origin of longitude, we have
\[
\theta' = \theta - n'u = \frac{1}{\sqrt{-1}}\log\frac{A - B}{A + B},
\]
($\sqrt{-1} = i$); and in consequence of this form of the expression for $\theta'$ ($= \frac{1}{i}\log$ into a logarithmic function), in passing to the trigonometrical functions $\sin\theta$, $\cos\theta$ the logarithm disappears altogether; and we have in a simple form the expressions for the actual functions $\sin\theta$, $\cos\theta$, through which $\theta$ enters into the formul\ae, and thus, Jacobi remarks, the barrier is cleared which stands in the way when the expression of an angle is reduced to an elliptic integral of the third class.

\textbf{198.} For the better expression of the results, Jacobi joins to the functions $H$, $\Theta$, considered in the ``Fundamenta Nova,'' the functions $\Theta_1 u = \Theta(K - u)$, $H_1 u = H(K - u)$, so that
\[
\sqrt{k}\sin\text{am}\,u = \frac{Hu}{i\sqrt{k}\,\Theta u}, \quad \sqrt{k}\cos\text{am}\,u = \frac{H_1 u}{\sqrt{k}\,\Theta u}, \quad \Delta\text{am}\,u = \frac{\Theta_1 u}{\Theta u};
\]
and then considering the cosine inclinations of the principal axes to the invariable axis and the revolving axes in the invariable plane, these are all fractions which, neglecting constant factors, have the common denominator $\Theta u$; $\alpha''$, $\beta''$, $\gamma''$ (as shown by Rueb's formula) have the numerators $H_1 u$, $Hu$, and $\Theta_1 u$ respectively; and $\alpha$, $\alpha'$ have the numerators $H(u + ia) \pm H(u - ia)$, $\beta$, $\beta'$ the numerators $H_1(u - ia) \pm H_1(u + ia)$, $\gamma$, $\gamma'$ the numerators $\Theta(u + ia) + \Theta(u - ia)$ respectively: there are also expressions of a similar form for the angular velocities about the axis of $x$ and $y$; that about the axis of $z$ (the invariable axis) having, as was known, the constant value $n$. The memoir is also very valuable analytically, as completing the systems of formul\ae\ given in the \textit{Fundamenta Nova} in reference to elliptic integrals of the third class.

\textbf{199.} It is worth noticing how the results connect themselves with Poinsot's theorem of the rolling and sliding cone: the velocity of the rolling motion depends only upon the position, on the cone, of the line of contact, so that the same series of velocities recur after any number of complete revolutions of the cone; that is, the total angle described by the line of contact in consequence of the rolling motion, consists of a part varying directly with the time (or say varying as $u$) and a periodic part; the former part combines with the similar term arising from the sliding motion; and the two together give Jacobi's term $n'u$.

\textbf{200.} Somoff's memoir (1851), written after Jacobi's Note in the \textit{Comptes Rendus}, but before the appearance of the memoir in Crelle, gives the demonstration of the greater part of Jacobi's results.

\textbf{201.} Booth's \textit{Theory of Elliptic Integrals} \&c.\ (1851) (contemporaneous with the publication of Poinsot's memoir of 1834) contains various interesting analytical developments, and, as an interpretation of them, the author obtains (p.~93) the theorem of the rolling and sliding cone. The investigations involve the elliptic integrals, but not the elliptic or Jacobian functions.

\textbf{202.} Richelot's two Notes (Crelle, tt.~XLII.\ and XLIV.) relate to the solution of the problem of rotation given in his memoir ``Eine neue L\"osung \&c.'' (1851). This is an application of Jacobi's theorem for the integration of a system of dynamical equations by means of the principal function $V$ (see my ``Report'' of 1857, art.~34). Retaining Richelot's letters $\phi$, $\psi$, $\theta$, which signify
\begin{itemize}[nosep]
    \item[$\phi$,] the longitude of the node,
    \item[$\theta$,] the inclination,
    \item[$\psi$,] the hour-angle,
\end{itemize}
the question is to find a complete solution of the partial differential equation
\[
\frac{dV}{dt} + \frac{1}{2A}\left(\frac{dV}{d\phi}\right)^2 + \frac{1}{2}\left(\frac{1}{A\sin^2\theta} + \frac{\cos^2\theta}{C}\right)\left(\frac{dV}{d\psi}\right)^2 + \frac{1}{2C}\left(\frac{dV}{d\theta}\right)^2 - \frac{\cos\theta}{A\sin^2\theta}\frac{dV}{d\phi}\frac{dV}{d\psi} = 0,
\]
that is, a solution involving (besides the constant attached to $V$ by a mere addition) three arbitrary constants; these are $a$, $b$, $p$. Writing in the first place $V = W + at + b\phi + p\psi$, the resulting equation for $W$ may be satisfied by taking $W$, a function of $\phi$ and $\theta$, without $\psi$ or $t$; and it is sufficient to have a solution involving only a single arbitrary constant.

This leads to a solution which is as follows:
\begin{align*}
V &= at + b\phi + p\psi \\
&\quad - \int\!\frac{d\theta}{\sqrt{\Theta}} \sqrt{(b - p\cos\theta)^2 - 2A(a - p^2/A - p^2/C)(\sin^2\theta)} \\
&\quad + \int\!\frac{d\theta}{\sqrt{\Theta}} \sqrt{(b - p\cos\theta)^2 - 2A(a - p^2/A - p^2/C)(\sin^2\theta) + \ldots}
\end{align*}
where $\Theta$ and the radicals are certain given functions of $a$, $b$, $p$, and of $\theta$ and $\phi$. The solution of the dynamical problem is then obtained by putting the differential coefficients $\frac{dV}{da}$, $\frac{dV}{db}$, $\frac{dV}{dp}$ equal to arbitrary constants $L$, $G$, $H$ respectively; the results are somewhat more simple than might be expected from the very complicated form of the function $V$. The foregoing results (although not by themselves very intelligible) will give an idea of the form in which the analytical solution in the first instance presents itself.

\textbf{203.} Richelot proceeds to remark that the solution in question, and the resulting integral equations of the problem, may be simplified in a peculiar manner by the method which he calls ``the integration by the spherical triangle.'' For this purpose he introduces a spherical triangle, the sides and angles whereof are $V$, $\lambda$, $\mu$; $N$, $\Lambda$, $M$, and then assuming $N$ constant, $M = \pi - \psi$,
\[
(B - C)\sin^2(\phi + \psi)\sin^2\Lambda + (A - B)\cos^2(\phi + \psi)\sin^2\Lambda = A - C,
\]
where $p$ and $q$ are constant, the solution is
\[
V = at - p(\psi - \lambda)\cos N - pM + p\int\!\cos\Lambda\,d(\phi + \psi);
\]
and that this expression leads to all the results almost without calculation.

\textbf{204.} I have quoted the foregoing results from the Note (Crelle, t.~XLII.), having seen, but without having studied, the Memoir itself: the results appear very interesting and valuable ones; but they seem to require a more complete geometrical development than they have received in the Memoir; and I am not able to bring them into connexion with the other researches on the subject.

\textbf{205.} The solution, \S~3 of Donkin's memoir ``On a Class of Differential Equations \&c.'' (part I.~1854), is given as an illustration of the general theory to which the memoir relates; it contains, however, some interesting geometrical developments in regard to the case ($A = B$) of two equal moments of inertia. I have not compared the results with those in my Note of 1849.

\vspace{1cm}
\noindent\rule{\textwidth}{0.5pt}
\vspace{0.5cm}

\noindent\textit{[End of Paper [298], Publication Pages 530---578]}

\end{document}
